% concatenated from Submitted_version__2_.tex
\documentclass[11pt]{article}
\usepackage{a4wide}

\usepackage[a4paper, margin=2.3cm]{geometry}
\usepackage{amsmath,amssymb,amsthm, comment, enumitem, array}
\usepackage{color}
\usepackage{parskip}
\usepackage{hyperref}

% Start TeXmacs macros
\usepackage{parskip}
\usepackage{setspace}
\usepackage{graphicx}
\usepackage{mathtools}
\usepackage[thinc]{esdiff}
\usepackage[euler]{textgreek}


% \usepackage[utf8]{vietnam}
\usepackage{babel}

\usepackage[capitalise, nameinlink]{cleveref}
\crefname{equation}{}{}
\crefname{figure}{{\sc Figure}}{{\sc Figure}}
\crefname{subsection}{Subsection}{Subsections}


\renewcommand{\baselinestretch}{1.2} 
\newtheorem{theorem}{Theorem}[section]
\newtheorem{proposition}[theorem]{Proposition}
\newtheorem{remark}[theorem]{Remark}
\newtheorem{corollary}[theorem]{Corollary}
\newtheorem{lemma}[theorem]{Lemma}
\newtheorem{question}[theorem]{Question}
\newtheorem{observation}[theorem]{Observation}
\newtheorem{definition}[theorem]{Definition}
\newtheorem{conjecture}[theorem]{Conjecture}
\newtheorem*{definition*}{Definition}
\newtheorem{exam}{Example}

\def\thang#1{\noindent
\textcolor{blue}
{\textsc{(Thang:}
\textsf{#1})}}
\def\thang#1{\noindent
\textcolor{blue}
{\textsc{(Thang:}
\textsf{#1})}}
\def\hung#1{\noindent
\textcolor{red}
{\textsc{(Hung:}
\textsf{#1})}}
\def\mq#1{\noindent
\textcolor{blue}
{\textsc{(MQ:}
\textsf{#1})}}
\newcommand{\hi}[1]{\textcolor{blue}{#1}}

\def\semin#1{\noindent
\textcolor{magenta}
{\textsc{(Semin:}
\textsf{#1})}}

\hypersetup{
    colorlinks=true,
    linkcolor = blue,
    citecolor=magenta,
    urlcolor=cyan,
}

%size of margin
% \setlength{\marginparwidth}{2.2cm}
% \setlength{\textwidth}{16cm} \setlength{\textheight}{20cm}
% \setlength{\oddsidemargin}{0.3cm} \setlength{\evensidemargin}{0.3cm}


\def\F{\mathcal{F}}
\def\M{\mathcal{M}}
\def\Fp{\mathbb{F}_p}
\def\Fbb{\mathbb{F}}
\def\Pbb{\mathbb{P}}
\def\Zbb{\mathbb{Z}}
\def\Rbb{\mathbb{R}}
\def\Cbb{\mathbb{C}}
\def\Qbb{\mathbb{Q}}
\def\Nbb{\mathbb{N}}
\def\R{\mathbb{R}}
\def\C{\mathcal{C}}
\def\A{\mathcal{A}}
\def\B{\mathcal{B}}
\def\E{E}\def\spt{\mathtt{spt}}
\def\RR{\mathcal{R}}
\def\S{\mathcal{S}}
\newcommand{\I}{\mathcal{I}}
\def \R{{\mathbb R}}
\def\rr{{\mathbb R}}
\def \F{{\mathbb F}}
\def \Z{{\mathbb Z}}
\def \1{\textbf{1}}
\def\supp{\hbox{supp\,}}

\DeclareMathOperator{\rank}{rank}
\DeclareMathOperator{\spn}{span}
%\def\F{\mathbb{F}_p}
\def\P{\mathcal{P}}
\def\L{\mathcal{L}}
\def\C{\mathcal{C}}
%\def\H{\mathcal{H}}
\def\S{\mathcal{S}}
\def\V{\mathcal{V}}
\def\Q{\mathcal{Q}}
\def\K{\mathcal{K}}
\def\U{\mathcal{U}}
\def\HH{\mathcal{H}}
\def\W{\mathcal{W}}
\def\A{\mathcal{A}}
\def\C{\mathcal{C}}
\def\B{\mathcal{B}}

\def\Ebf{\textbf{E}}
\def\SL{SL}


\def\l({\left(}
\def\r){\right)}

\def\lu{\left[}
\def\ru{\right]}

\def\lv{\left\vert}
\def\rv{\right\vert}

\def\lo{\left\{}
\def\ro{\right\}}

\def\lV{\left\Vert}
\def\rV{\right\Vert}

\def\la{\left\langle}
\def\ra{\right\rangle}

\def\lce{\left\lceil}
\def\rce{\right\rceil}

\def\lf{\left\lfloor}
\def\rf{\right\rfloor}

\def\rar{\rightarrow}
\def\lar{\leftarrow}

\newcommand{\sss}{\mathcal S}
\newcommand{\kkk}{\mathcal K}
\newcommand{\uuu}{\mathcal E}



\begin{document}
\title{Packing sets under finite groups via algebraic incidence structures}
\author{Norbert Hegyv\'{a}ri\thanks{E\"otv\"os University and associated member of Alfréd Rényi Institute, Hungary. Email: hegyvari@renyi.hu}\and Le Quang Hung\thanks{Hanoi University of Science and Technology, Vietnam. Email: hung.lequang@hust.edu.vn}\and Alex Iosevich\thanks{Department of Mathematics, University of Rochester, USA. Email: iosevich@math.rochester.edu}\and Thang Pham\thanks{Institute for Mathematics and Interdisciplinary Sciences, Xidian University. \newline
\hspace*{0.5cm} Email: {\tt thangpham.math@gmail.com}}}
\date{}
\maketitle

\begin{abstract}
Let $G$ be a finite group acting on a vector space $V = \mathbb{F}_p^n$ over a prime field. Given finite sets $S \subset G$ and $E \subset V$, we study the restricted orbit union $S(E) = \bigcup_{g\in S} g(E)$ and establish quantitative lower bounds for $|S(E)|$ in terms of $|S|$, $|E|$, and natural structural conditions. This finite field packing problem has connections to distance geometry, configuration counting, and expanding graphs. For $G = SL_2(\mathbb{F}_p)$ acting on $\mathbb{F}_p^2$, we prove that $$|S(E)| \gg \min\left\lbrace p^2, \frac{|S||E|}{p^2}\right\rbrace,$$ which is sharp. Under geometric non-concentration conditions on $E$ and subgroup-avoidance hypotheses on $S$, we obtain a power-saving improvement of the form $$|S(E)|\gg \min \left\lbrace p^2, ~\max\left\lbrace\frac{|S||E|}{pk}, ~\frac{|S|^{\frac{1}{2}}|E|}{p^{\frac{1-\epsilon}{2}}k^{\frac{1}{2}}}\right\rbrace \right\rbrace,$$ where $k$ bounds the radial multiplicity of $E$. For small sets $|E| \leq p$, we establish optimal bounds using weighted incidence theory. Analogous results are proved for the first Heisenberg group $\mathbb{H}_1(\mathbb{F}_p)$ acting on $\mathbb{F}_p^3$. Our approach reformulates the problem as an incidence question in a bipartite action graph. The proofs combine Fourier analytic techniques, energy estimates, point-line incidence bounds, and area-energy inequalities for skew dot products. The methods extend classical sum-product type problems and incidence theory to noncommutative group actions.
\end{abstract}



\textbf{Keywords}: Packing sets, incidence geometry, expanding maps, finite fields, group actions

\textbf{MSC-2020}: 11T30, 51B05

\tableofcontents
\section{Introduction}

Let $E$ be a Borel set in $\mathbb{R}^n$ and let $S$ be a set of maps from $\mathbb{R}^n$ to $\mathbb{R}^n$. We define 
\begin{align*}
    S(E):=\bigcup_{f\in S}f(E)=\{f(y): f\in S, y\in E\}\subset \R^n.
\end{align*}
The packing problem asks if it is possible for a set of zero $n$-dimensional Lebesgue measure to contain the image $f(E)$ for all $f\in S$. The study of this problem has a well-known history in the literature, for example, there are sets in the plane of zero Lebesgue measure containing  a line segment of unit length in every direction (see \cite{Besicovitch20} and \cite{Besicovitch28}), a circle of radius $r$ for all $r>0$ (see \cite{BesicovitchRado68} and \cite{Kinney68}), or a circle centered at $x$ for all $x$ on a given straight line (see \cite{Talagrand80}). In another direction, the question of finding conditions on $S$ and $E$ such that $S(E)$ has positive Lebesgue measure has also received a lot of attention. Bourgain \cite{Bourgain85} and independently Marstrand \cite{Marstrand87} proved that given a set of circles in the plane, if the centers form a set of positive Lebesgue measure, then the union of circles also has positive Lebesgue measure. Wolff \cite{Wolff97} strengthened this result by showing that if the set of centers has Hausdorff dimension $s$, $0<s\leq 1$, then the dimension of the set of union of circles is at least $1+s$. For the most recent progress, we refer the reader to \cite{Iosevichetal} and references therein. 

Let $\mathbb{F}_p$ be a prime field of order $p$, where $p$ is an odd prime. Let $G$ be a finite group acting on $V=\mathbb{F}_p^n$.
Given finite sets $S\subset G$ and $E\subset V$, define the restricted orbit union set
\[
S(E):=\{gy:\ g\in S,\ y\in E\}=\bigcup_{g\in S} g(E).
\]
In this setting, the packing question translates into quantitative lower bounds for the size of $S(E)$. This paper provides such bounds for $|S(E)|$ in terms of $|S|$, $|E|$, and natural non-concentration conditions.


Beyond its intrinsic expanding flavor, it is worth mentioning that restricted orbit union expansion implies several configuration-counting consequences, including pinned triangle congruence and distance questions. We briefly indicate this connection as follows. When $S$ is a set of rigid-motions, i.e. a set of maps preserving distances, the author of \cite{Pham2024} proved a lower bound on the number of congruence classes of triangles with one fixed side-length in the plane. By setting $S=\{g\}\times E'$, for some rotation $g\in O(n)$ and $E'\subset \mathbb{F}_p^n$, and $S(E)=\{gy-x\colon y\in E, x\in E'\}$, the lower bound $|S(E)|\gg p^n$ implies that the distance set between $gE$ and $E'$, denoted by $\Delta(gE, E')$, covers a positive proportion of all distances in $\mathbb{F}_p$. The authors of \cite{PS} proved that under the conditions $|E|, |E'|\ge p$, for almost every rotation $g$ in the plane $\mathbb{F}_p^2$, $|gE-E'|\gg p^2$, and obtained $|\Delta(gE, E')|\gg p$. Note that the Erd\H{o}s-Falconer distance conjecture in two dimensions corresponds to the case $g=id$. 

%Our main contribution is an incidence-theoretic framework for restricted orbit unions.
%We encode the action relation $g y=x$ as a bipartite action-incidence graph and prove sharp incidence bounds. In the cases $G=SL_2(\mathbb{F}_p)$ acting on $\mathbb{F}_p^2$ and $G=\mathbb{H}_1(\mathbb{F}_p)$ acting on $\mathbb{F}_p^3$,
%this implies optimal one-step expansion bounds for $|S(E)|$, together with $\epsilon$-improvements
%under subgroup/fiber non-concentration and geometric non-concentration.
 




\paragraph{Notations:} Throughout the paper, we will write $X \ll_\alpha Y$ if $X \le CY$, where $C>0$ is a constant depending on $\alpha$. If it is clear from the context what $C$ should depend on, we may write only $X \ll Y$. If $X \ll Y $ and $Y\ll X$, we write $X \sim Y$. Furthermore, $X \lesssim Y$ if $X \ll \log_2p\cdot Y$.  

\subsection{The special linear group \texorpdfstring{$G=SL_2(\mathbb{F}_p)$}{SL2(Fp)}}
We first start with some observations. 


    
{\bf Observation 1:} Since each $\theta\in SL_2(\mathbb{F}_p)$ acts bijectively on $\mathbb{F}_p^2$, we have $|S(E)|\ge |E|$. One might hope that $|S(E)|\ge |S|$, which implies $|S(E)|\ge |S|^{\frac{1}{2}}|E|^{\frac{1}{2}}$. However, the estimate $|S(E)|\ge |S|$ might not be true. For example, take $E=\{(0, 1)\}$ and $S$ being a set of matrices $\theta$ in $SL_2(\mathbb{F}_p)$ such that $\theta (0, 1)=(1, 0)$, then, by Lemma \ref{lmmay1}, $S$ can be as large as $\sim p$, and in this case, one has $|S(E)|=1$. 



There are sets $S$ and $E$ such that $|S(E)|=|S|^{\frac{1}{2}}|E|^{\frac{1}{2}}$. Indeed, let $A\subset B $ be two subgroups of $\Fbb_p^*$. For each $[x] \in B/A$, fix $x' \in [x]$. Let $S_{[x]}$ be a subset of $\lo \theta \in \SL_2 \l( \Fbb_p \r) \colon \theta (0,1) =(0,x') \ro$ such that $\lv S_{[x]} \rv = \lv B\rv$. Let $E= \lo  (0,a) \colon a\in A \ro$. Then, we have
    \[ S_{[x]} (E) = \{ 0\} \times [x] .\]
    Therefore, if we choose $S = \bigcup_{[x]\in B/A} S_{[x]}  $, we have $\lv S \rv = \lv B\rv\lv B/A\rv$ and
    \[ S (E) = \bigcup_{[x]\in B/A} \{ 0\} \times [x] =\{ 0\} \times B. \]
    Then $\lv S (E)\rv = \lv B\rv = (\lv B\rv \lv B/A\rv)^{\frac{1}{2}} \lv A\rv^{\frac{1}{2}} = \lv S \rv^{\frac{1}{2}} \lv E \rv^{\frac{1}{2}} $.

{\bf Observation 2:} There are sets $S$ and $E$ such that $\lv S(E) \rv \sim \frac{\lv S\rv \lv E \rv}{p^2}>|E|$. Indeed, fix $0<\epsilon<1$, let $E$ be the set of points on the line $\{y=0\}$. Let $S \subset \SL_2 \l( \Fbb_p \r)$ be a set of all matrices $\theta$ such that $\theta (1,0) \in E'$, where $E'$ is a set of $p^{1+\epsilon}$ points on $p^{\epsilon}$ lines passing through the origin. Then, by Lemma \ref{lmmay1}, we have $\lv S\rv \sim p^{2+\epsilon}$. Therefore,
\[ \lv S(E)\rv =p^{1+\epsilon} \sim \frac{p^{2+\epsilon}\cdot p}{p^2} \sim \frac{\lv S\rv \lv E\rv}{p^2}. \]



In the first result, we prove that the lower bound $|S||E|/p^2$ actually holds for all sets $S$ and $E$. 
\begin{theorem}\label{thm2}
    Let $E\subset \mathbb{F}_p^2$ and $S\subset SL_2(\mathbb{F}_p)$. We have 
    \[|S(E)|\gg \min \left\lbrace p^2, ~\frac{|S||E|}{p^2} \right\rbrace.\]
\end{theorem}
It follows from this theorem that if $|S||E|\gg p^4$ then $|S(E)|\gg p^2$. This condition is optimal in the sense that for all $\epsilon>0$ there exist sets $S$ and $E$ with $|S||E|\gg p^{4-\epsilon}$ such that $|S(E)|=o(p^2)$. Indeed, for $\epsilon>0$, by choosing $p$ large enough, we can find a cyclic subgroup $A$ of $\mathbb{F}_p^*$ with $|A|\sim p^{1-\epsilon}$. Let $S$ be the set of matrices $\theta$ in $SL_2(\mathbb{F}_p)$ such that $\theta(0, 1)\in \{(-x, 0)\colon x\in A \}$. Each such matrix is of the form
\[\begin{bmatrix}
    *&-x\\
    x^{-1} &0
\end{bmatrix},\]
and $|S|\sim p^{2-\epsilon}$. We now let $E$ to be the set of points of the form $(y, *)$ with $y\in A$ and $*\in \mathbb{F}_p$. Then $|E|\sim p^{2-\epsilon}$. Moreover, $S(E)$ is covered by the lines of the from $y=\lambda$ with $\lambda\in A$. So $|S(E)|\ll p^{2-\epsilon}$.

Under subgroup non-concentration for $S$ and geometric non-concentration for $E$, the next theorem presents quantitatively stronger expansion bounds than those in Theorem~\ref{thm2}.



\begin{theorem}\label{thm555}
For any $\gamma\in (0, 1)$, there exists $\epsilon=\epsilon(\gamma)>0$ such that the following holds. 
     Let $E\subset \mathbb{F}_p^2$ and $S\subset SL_2(\mathbb{F}_p)$. Assume $p$ is a sufficiently large prime, $S$ is a symmetric subset of $\SL_2 (\Fbb_p )$ such that $p^{\gamma} < |S| < p^{3-2\gamma } $ and $ |S \cap gH| < p^{-\frac{\gamma}{2}}|S|$ for any subgroup $H \subsetneq \SL_2 (\Fbb_p) $ and $g\in \SL_2(\Fbb_p)$, and any line through the origin contains at most $k$ points from $E$, then 
     \[|S(E)|\gg \min \left\lbrace p^2, ~\max\left\lbrace\frac{|S||E|}{pk}, ~\frac{|S|^{\frac{1}{2}}|E|}{p^{\frac{1-\epsilon}{2}}k^{\frac{1}{2}}}\right\rbrace \right\rbrace.\]
%Moreover, If  $|E|\ge 4p$ and $|S|\gg p^2$, then there exists $x\in E$ such that $|S(E-x)|\gg p^2$.
\end{theorem}


\paragraph{Sharpness:} We observe that conditions on $S$ in the above theorem are natural for improvements. Indeed, let $\ell_1$  and $\ell_2$ be two lines passing through the origin and $S$ be the set of all $\theta \in \SL_2(\Fbb_p)$ such that $\theta (\ell_1) =\ell_2$. Then, let $E= \ell_1$, we have $S(E)=\ell_2$, and by Lemma \ref{lmmay1}, $|S|\sim p^2$, so
    \[ |S(E)| = p \sim \frac{|S|^{\frac{1}{2}}|E|}{p}. \]
    Moreover, we have $S=gH$, where $H$ is the group of matrices $\theta \in \SL_2(\Fbb_p)$ such that $\theta (\ell_1) =\ell_1$ and $g\in \SL_2(\Fbb_p)$ is a matrix such that $g(\ell_1) =\ell_2$. 
 
We now turn our attention to the case of small sets. If $S$ and $E$ are arbitrary sets of small size, then based on Observation 1 and Observation 2, one can guess of the existence of sets such that the size of $S(E)$ is the same as the trivial lower bound $|E|$. For example, let $A$ be a subgroup of $\mathbb{F}_p^*$, and let $E=\{\lambda (1, 0)\colon \lambda\in A\}$, and $S$ be a set of matrices 
$\theta$ such that $\theta$ maps at least one point in $E$ to $(0, 1)$. Depending on the form of $p$, one can choose $E$ of arbitrary small size, and the size of $S$ can also be chosen arbitrary large in the range $(0, p|E|)$ such that $|S(E)|=|E|$. 

In the setting of small sets, we prove the following optimal result. 
\begin{theorem}\label{thm345} 
  Let $S \subseteq \SL_2 \l( \Fbb_p \r)$, and $E \subseteq \Fbb_p^2\setminus \lo (0, 0)\ro$ be such that $|E| \le p$, any line passing through the origin contains at most $k_1$ points from $E$, and $E$ determines at most $k_2$ distinct directions through the origin. Then,
  \[ \lv S(E)\rv \gtrsim \min \lo  \frac{\lv E\rv \lv S\rv^{\frac{1}{2}}}{k_1^{\frac{1}{2}}k_2^{\frac{1}{2}}}, \frac{\lv E\rv^{\frac{1}{2}} \lv S\rv^{\frac{1}{2}}}{k_1^{1/4}}, \frac{\lv E\rv \lv S\rv^{\frac{1}{2}}}{k_1}, \frac{\lv E\rv^2}{k_1} \ro .\]
Here, by ``directions" we mean the number of lines passing through the origin that are required to cover the whole set $E$.
\end{theorem}
In order to have $|S(E)|> |E|$, one would need the conditions that 
\[|E|\gg k_1, ~~|S|\ge \max \left\lbrace k_1^2, k_1^{1/2}|E|, k_1k_2  \right\rbrace.\]


The lower bound of Theorem \ref{thm345} is attainable from the following example.

For a constant $0<c<1$, let $\lo \ell_i \ro_{i=1}^{\frac{p}{c}}$ be a family of $\frac{p}{c}$ lines passing through the origin. For each $1\le i \le \frac{p}{c}$, choose $a_i \in \ell_i \setminus \lo (0,0) \ro$, and set $E = \lo a_i \colon 1 \le i \le \frac{p}{c} \ro$, $S = \SL_2(\Fbb_p)$. Then we have $k_2=|E|=\frac{p}{c} ,k_1=1$, and $|S(E)|=p^2-1$. 
    Thus, 
     \[ \lv S(E)\rv \sim  \min \lo  \frac{\lv E\rv \lv S\rv^{\frac{1}{2}}}{k_1^{\frac{1}{2}}k_2^{\frac{1}{2}}}, \frac{\lv E\rv^{\frac{1}{2}} \lv S\rv^{\frac{1}{2}}}{k_1^{1/4}}, \frac{\lv E\rv \lv S\rv^{\frac{1}{2}}}{k_1}, \frac{\lv E\rv^2}{k_1} \ro .\]

As in Theorem \ref{thm555}, we now provide an $\epsilon$-improvement under non-concentration conditions.

\begin{theorem}\label{thm3'}
For any $\gamma\in (0, 1)$, there exists $\epsilon=\epsilon(\gamma)>0$ such that the following holds. 
  Let $p$ be a sufficiently large prime, and $E \subseteq \Fbb_p^2\setminus \lo (0, 0)\ro$ such that $|E| \le p$, any line passing through the origin contains at most $k_1$ points from $E$ and $E$ determines at most $k_2$ distinct directions through the origin. Let $S \subseteq \SL_2 \l( \Fbb_p \r)$ be a symmetric subset such that $p^{\gamma} < \lv S\rv < p^{3-2\gamma} $ and $\lv S \cap gH\rv < p^{-\frac{\gamma}{2}} \lv S\rv $ for any subgroup $H \subsetneq \SL_2 \l( \Fbb_p \r)$ and $g \in \SL_2 \l( \Fbb_p \r)$. Then,
  \[ \lv S(E)\rv \gtrsim \min \lo  \frac{\lv E\rv \lv S\rv^{\frac{1}{2}}p^{\epsilon /2}}{k_1^{\frac{1}{2}}k_2^{\frac{1}{2}}}, \frac{\lv E\rv^{\frac{1}{2}} \lv S\rv^{\frac{1}{2}}p^{\epsilon /2}}{k_1^{1/4}}, \frac{\lv E\rv \lv S\rv^{\frac{1}{2}}p^{\epsilon /2}}{k_1}, \frac{\lv E\rv^2}{k_1} \ro .\]
\end{theorem}

Our approach to quantitative lower bounds for $|S(E)|$ is to recast the orbit-union problem as an incidence problem attached to the group action.
Instead of working directly with
\[
S(E)=\{gy:\ g\in S,\ y\in E\},
\]
we encode the relation $gy=x$ as incidences in a bipartite \emph{action-incidence graph} whose left vertices are pairs $(x,y)$ and whose right vertices are group elements $g\in S$.
This reformulation reduces expansion estimates for $|S(E)|$ to proving \emph{incidence discrepancy} bounds, i.e. upper bounds for
\[
\Bigl|I(A\times B,S)-\tfrac{|A||B||S|}{p^{2}}\Bigr|
\]
in terms of natural structural parameters of $A,B$, and $S$.

For $G=\SL_2(\F_p)$ acting on $\F_p^2$, we first establish a sharp universal estimate of the form
\[
\Bigl| I(A\times B,S)-\tfrac{|A||B||S|}{p^{2}} \Bigr|
\ \ll\ p\sqrt{|A||B||S|}+|S|
\]
(Theorem~\ref{incidence}).
The proof is Fourier-analytic: after expanding the incidence constraint and applying Cauchy--Schwarz, the key step is a geometric reduction using the fact that $\SL_2(\F_p)$ preserves the skew dot product (equivalently, signed areas), cf.\ Lemma~\ref{lmmay2}.
This converts the main contribution to a bound for the ``area-energy'' quantity
\[
\#\{(x_1,x_2,y_1,y_2)\in A^2\times B^2:\ x_1\cdot x_2^\perp=y_1\cdot y_2^\perp\},
\]
which is controlled by known $L^2$ estimates for the skew dot-product/determinant function.
Applying the resulting incidence bound with $A=S(E)$ and $B=E$, and comparing with the trivial lower bound
$I(S(E)\times E,S)\ge |S||E|$, implies the one-step expansion bound in Theorem~\ref{thm2}.

To obtain stronger bounds when additional structure is available, we refine the incidence analysis via an \emph{energy decomposition}.
Under geometric non-concentration on $E$ (at most $k$ points on any line through the origin) and subgroup non-concentration for $S$, we prove an $\epsilon$-improved incidence estimate (Theorem~\ref{thm: 4.12})
\[
I(S(E)\times E,S)\ \ll\ \frac{|S(E)|^{1/2}|E||S|}{p}
\ +\ k^{\frac{1}{4}}p^{\frac{1-\epsilon}{4}}|S(E)|^{\frac{1}{2}}|E|^{\frac{1}{2}}|S|^{\frac{3}{4}}.
\]
The mechanism is a Cauchy--Schwarz reduction to a collision count, which naturally introduces two energies:
(i) a geometric area-energy of $E$, controlled in terms of $k$, and
(ii) the multiplicative energy $E(S,S)=\#\{(a,b,c,d)\in S^4:\ ab=cd\}$.
While $E(S,S)$ can be as large as $\sim |S|^{3}$ in general, Bourgain--Gamburd $L^2$-flattening yields the power-saving bound
$E(S,S)\ll |S|^{3}p^{-\epsilon}$ under subgroup avoidance \cite{B.A.08}.
Incorporating this saving into the refined incidence estimate produces the $\epsilon$-improved expansion statements in Theorems~\ref{thm555} and~\ref{thm3'}.

Finally, in the small-set regime $|E|\le p$, the Fourier methods become ineffective and we instead use incidence geometry.
We prove a weighted/multi-set extension of the Stevens--de Zeeuw point-line incidence theorem (Theorem~\ref{eqn:IPLMS}), obtained via a dyadic decomposition of multiplicities.
This weighted incidence input is then used to bound skew dot-product energies for small configurations, leading to sharp expansion estimates in terms of the \emph{radial multiplicity} $k_1$ (maximum number of points of $E$ on a line through the origin) and the \emph{number of directions} $k_2$ determined by $E$, as in Theorems~\ref{thm345} and~\ref{thm3'}.





One might ask about the higher dimensional case. Let's assume $n=3$ for simplicity. It is not hard to check that if we have two triples $(x, y, z)$ and $(x', y', z')$ such that each forms an independent system, then there exists unique $\theta\in SL_3(\mathbb{F}_p)$ such that $\theta x=x'$, $\theta y=y'$, and $\theta z=z'$ if and only if $\det(x, y, z)=\det(x', y', z')$. Here $\det(x, y, z)$ is the determinant of the matrix with columns $x$, $y$, and $z$. If one wishes to apply the method in the two dimensions, then an estimate on the number of tuples $(x, y, z, x', y', z')\in E^6$, $E\subset \mathbb{F}_p^3$, such that $\det(x, y, z)=\det(x', y', z')$ is needed in the first step. This is the $L^2$-norm version of earlier results studied in \cite{covert, lavinh}. Notice that this framework only solves the case of large sets, and for the case of small sets in higher dimensions, it appears to be a hard problem due to the limited understanding of incidence bounds. We plan to address this in a subsequent paper.

\subsection{The first Heisenberg group 
\texorpdfstring{$G=\mathbb{H}_1(\mathbb{F}_p)$}{H1(Fp)}}

We now move to the case of the Heisenberg group. Let $\mathbb{F}_p$ be a prime field, we denote by $\mathbb{H}_1(\mathbb{F}_p)$ the first Heisenberg group over $\mathbb{F}_p$, i.e. the group of matrices of the form 
\[ [x, y, t]:=\begin{pmatrix}1&x&t\\
0&1&y\\
0&0&1\end{pmatrix}, ~x, y, t\in \mathbb{F}_p.\]
In the group $\mathbb{H}_1(\Fbb_p)$, we have 
\[[x, y, t]\cdot [x', y', t']=\left(x+x', y+y', t+t'+\frac{xy'-yx'}{2}\right).\]

This section uses the notation $X(E)$ in order to distinguish with the case of $SL_2(\mathbb{F}_p)$.

We first look at the following observations. 

{\bf Observation 3:} 
Let $X=\mathbb{H}_1(\Fbb_p)$ and let $E$ be the set of all points with the third coordinate belonging to a set of size $\alpha p$ in $\mathbb{F}_p$. A direct calculation gives $|X(E)|\le \alpha p^3$. This example tells us that in order to cover the whole space or a positive proportion of all elements in $\mathbb{F}_p^3$, the condition $|E|\gg p^3$ is needed. 

{\bf Observation 4:} There are sets $X$ and $E$ of arbitrary large such that $|X(E)|\ll |X|$. Let $A\subset\mathbb{F}_p$ be an arbitrary set, let $X$ be the set of matrices $[a,b,c]\in \mathbb{H}_1(\Fbb_p)$ with $a, c\in \mathbb{F}_p$ and $b\in A$, let $E$ be the set of points $(x, y, z)$ with $x\in \mathbb{F}_p,~ y\in A,$ and $z\in A$. Then we have $|X|=p^2|A|$ and $|E|=p|A|^2$. A direct computation shows that $|X(E)|\ll p^2|A|= |X|$. Thus, for any $0<\epsilon<2$, there exist sets $E\subset \mathbb{F}_p^3$ and $X\subset \mathbb{H}_1(\mathbb{F}_p)$ with $|E|, |X|\ge p^{3-2\epsilon}$ such that $|X(E)|\ll |X|$.

Let $\pi_{23}\colon \mathbb{F}_p^3\to \mathbb{F}_p^2$ defined by $\pi_{23}(x, y, z)=(y, z)$. The following example shows that if the pre-image set $\pi_{23}^{-1}(y, z)\cap E$ is of large size for all $(y, z)\in \pi_{23}(E)$, then the trivial bound $|E|$ might be best possible.

{\bf Observation 5:} There are sets $X$ and $E$ such that $|X(E)|\ll |E|$. Let $A\subset\mathbb{F}_p$ be an arithmetic progression. Let $E$ be a set in $\mathbb{F}_p^3$ such that each point in $E$ has the last two coordinates belonging to $A$. Assume that for each $(y, z)\in \pi_{23}(E)$, we have $|\pi_{23}^{-1}(y, z)\cap E|\sim p$. Then, for all sets $X$ containing of matrices of the form $[*, 1, *]\in \mathbb{H}_1(\Fbb_p)$, we have $|X(E)|\ll |E|$. 

Under a non-concentration condition, the main theorem in this section reads as follows. 

\begin{theorem}\label{main-H}
Let $\epsilon\ge 0$, $X$ be a subset of $\mathbb{H}_1(\Fbb_p)$, and $E$ be a set in $\mathbb{F}_p^3\setminus (\mathbb{F}_p^2\times \{0\})$. Assume that for each $(y, z)\in \mathbb{F}_p^2$, we have $|\pi_{23}^{-1}(y, z)\cap E|\le p^{1-\epsilon}$, then we have
\[|X(E)|\gg \min \left\lbrace p^3, ~\frac{|X||E|}{p^{3-\frac{\epsilon}{2}}} \right\rbrace.\]
\end{theorem}

Note that when $\epsilon=0$, it implies directly that 
\[|X(E)|\gg \min \left\lbrace p^3, ~\frac{|X||E|}{p^{3}} \right\rbrace,\]
which is sharp by Observation 3. 

As in the $\SL_2(\F_p)$ setting, we prove Theorem~\ref{main-H} by encoding the action relation
$\theta y=x$ as incidences in an associated bipartite action-incidence graph. However, the
underlying invariant structure is different.  Writing
\[
[a,b,c]:=\begin{pmatrix}1&a&c\\ 0&1&b\\ 0&0&1\end{pmatrix}\in \mathbb{H}_1(\F_p),
\qquad
[a,b,c](x,y,z)=(x+ay+cz,\; y+bz,\; z),
\]
we see that the third coordinate is preserved; accordingly we restrict to points with $z\neq 0$
(and likewise in Theorem~\ref{H-inci}).

The key step in the incidence analysis is to understand when a \emph{single} group element
$[a,b,c]$ can simultaneously map two points $(x,y,z),(u,v,w)$ to two targets
$(x',y',z'),(u',v',w')$.  From the explicit action one immediately gets $z'=z$, $w'=w$ and
\[
b=\frac{y'-y}{z}=\frac{v'-v}{w},
\]
so compatibility forces the bilinear constraint
\begin{equation}\label{bracket}
z\,(v'-v)=w\,(y'-y)
\qquad
\text{(equivalently } y\,w+z\,v'=y'\,w+v\,z\text{)},
\end{equation}
together with $z=z'$ and $w=w'$.  After a Fourier expansion of the incidence relation and an
application of Cauchy--Schwarz, the discrepancy reduces to bounding a \emph{weighted} count of
quadruples satisfying this bracket equation, where the weights are precisely the multiplicities of
fibers of the projection $\pi_{23}(x,y,z)=(y,z)$.  A weighted orthogonality estimate
(Lemma~\ref{weight-H}) combined with the fiber non-concentration hypothesis
$|\pi_{23}^{-1}(y,z)\cap B|\le p^{1-\epsilon}$ gives
\[
\Bigl|I(A\times B,X)-\tfrac{|A||B||X|}{p^{3}}\Bigr|
\ \ll\ p^{\frac{3-\epsilon}{2}}\,|A|^{\frac{1}{2}}|B|^{\frac{1}{2}}|X|^{\frac{1}{2}},
\]
i.e. Theorem~\ref{H-inci}, with degenerate configurations handled by a
separate analysis.
Applying the incidence bound with $A=X(E)$ and $B=E$ then gives Theorem~\ref{main-H}.
Finally, Observation~5 shows that fiber control is essentially necessary: if many points of $E$
lie in single $\pi_{23}$-fibers (of size $\sim p$), even large $X$ may fail to expand beyond the
trivial scale.  We therefore do not pursue small-set analogues in $\mathbb{H}_1(\F_p)$ here, as the
corresponding three-dimensional constraint of (\ref{bracket}) lead to substantially more intricate incidence
geometry.



%One might also ask about a version of Theorem \ref{main-H} for small sets by using incidence bounds. However, we find it too complicated to pursue in this direction. For simplicity of this paper, we pose it as an open question. 

\section{Incidence structures spanned by \texorpdfstring{$SL_2(\mathbb{F}_p)$}{SL2(Fp)}}
Let $x,y \in \Fbb_p^2$ and $\theta \in \SL_2 \l( \Fbb_p \r) $, we say $(x,y)$ is incident to $\theta$ if $\theta y=x$. Let $P=A\times B \subseteq \Fbb_p^2 \times \Fbb_p^2$ and $S\subseteq \SL_2 \l( \Fbb_p \r) $, we denote the number of incidences between $P$ and $S$ by $I(P,S)$.

In this section, we prove the following incidence bound. 

\begin{theorem}\label{incidence} Let $P=A\times B \subseteq \mathbb{F}_p^2\times \mathbb{F}_p^2$ and $S \subseteq \SL_2 \l( \Fbb_p \r) $. Then, we have
    \[\lv I(P,S)- \frac{\lv P\rv \lv S\rv}{p^2} \rv \ll p\sqrt{\lv S\rv \lv P\rv}+|S|.\]
% \begin{itemize}
%     \item[(1)] if $|P||S|\ge p^6$, we have 
%     \[ I(P,S) \ll \frac{|P||S|}{p^2}, \]
%     \item[(2)] if $|P||S| < p^6$, we have 
%     \[ I(P,S) \ll q\sqrt{|P||S|} + |S|. \]
% \end{itemize}
Moreover, let $k_A$ and  $k_B$ be the maximal number of points from $A$ and $B$ on a line passing through the origin, respectively. Assume that $\min\{k_A, k_B\}= k$,
then
 \[\lv I(P,S)- \frac{\lv P\rv \lv S\rv}{p^2} \rv \ll p^{\frac{1}{2}}k^{\frac{1}{2}}\sqrt{\lv S\rv \lv P\rv}+|S|.\]
\end{theorem}

We have some comments on this theorem. 
\begin{enumerate}
    \item This theorem is sharp, i.e. there are sets $P$ and $S$ such that the upper bound is attained. 

     Let $P= A \times B$ where $A,B$ are the sets of points on two lines through the origin. Let $S$ be the set of all matrices $\theta$ in $\SL_2 \l( \Fbb_p \r)$ such that $\theta (B)=A$. Then, we have $\lv S\rv =p(p-1), \lv P\rv =p^2$, and
    \[ I(P,S) = p(p-1) + p(p-1)^2.  \]
    Hence,
    \[ \lv I(P,S) - \frac{\lv P\rv \lv S\rv}{p^2} \rv = p(p-1)^2 \sim p\sqrt{\lv P\rv \lv S\rv} + \lv S\rv .\]
    
    \item The same result holds for general sets $P\subset \mathbb{F}_p^4$ instead of sets of Cartesian product structures.
\end{enumerate}
If we assume some structural conditions on $S$, the upper bound of the above theorem can be improved further. 

\begin{theorem}\label{thm: 4.12}
For $\gamma\in (0, 1)$, there exists $\epsilon=\epsilon(\gamma)>0$ such that the following holds.
    Let $p$ be a sufficiently large prime. Let $P=A\times B \subseteq (\Fbb_p^2 \times \Fbb_p^2 ) \setminus \lo (0, 0, 0, 0)\ro$, and let $S$ be a symmetric subset of $\SL_2 (\Fbb_p )$ such that $p^\gamma < |S| < p^{3-2\gamma } $ and $ |S \cap gH| < p^{\frac{-\gamma}{2}}|S|$ for any subgroup $H \subsetneq \SL_2 (\Fbb_p) $ and $g\in \SL_2(\Fbb_p)$. Then, we have 
    \[ I(P,S) \ll \frac{|A|^{\frac{1}{2}}|B||S|}{p} + p^{\frac{2-\epsilon}{4}}|A|^{\frac{1}{2}}|B|^{\frac{1}{2}}|S|^{\frac{3}{4}}.\] Moreover,  if any line passing through the origin contains at most $k$ points from $B$, we have
    \[ I(P,S) \ll \frac{|A|^{\frac{1}{2}}|B||S|}{p} + k^{\frac{1}{4}}p^{\frac{1-\epsilon}{4}}|A|^{\frac{1}{2}}|B|^{\frac{1}{2}}|S|^{\frac{3}{4}}.\]
\end{theorem}

%\thang{Hung, if you do not have any conditions on $S$, then is it correct that the upper bounds 
% \[ I(P,S) \ll \frac{|A|^{\frac{1}{2}}|B||S|}{p} + p^{\frac{2}{4}}|A|^{\frac{1}{2}}|B|^{\frac{1}{2}}|S|^{\frac{3}{4}}  ,\]
%    and
%    \[ I(P,S) \ll \frac{|A|^{\frac{1}{2}}|B||S|}{p} + k^{\frac{1}{4}}p^{\frac{1}{4}}|A|^{\frac{1}{2}}|B|^{\frac{1}{2}}|S|^{\frac{3}{4}}  ,\]
% are sharp?
%} 
The following example shows that conditions on $S$ are necessary to obtain non-trivial improvements.

Let $\ell_1$ and $\ell_2$ be two lines passing through the origin. Let $A, B$ be the set of points on $\ell_1, \ell_2$, respectively. Let $S$ be the set of all $\theta\in SL_2(\mathbb{F}_p)$ such that $\theta(\ell_2)=\ell_1$, and $H$ be the group of matrices $\theta\in SL_2(\mathbb{F}_p)$ such that $\theta(\ell_2)=\ell_2$. Then 
\[ I(P,S) \sim p|A||B| \sim p^3  \sim \frac{|A|^{\frac{1}{2}}|B||S|}{p} + p^{\frac{2}{4}}|A|^{\frac{1}{2}}|B|^{\frac{1}{2}}|S|^{\frac{3}{4}}   . \]
Moreover, it is not hard to check that $S=gH$ for some $g\in \SL_2(\Fbb_p)$.



%is a matrix such that $g(B)=A$ and $H$ is stabilizer subgroup of $B$ under $\SL_2(\Fbb_p)$. Indeed, for any $\theta \in S$, we have $g^{-1}\circ \theta (B)= B$, so $g^{-1}\theta \in H$ or $\theta \in gH$. Thus, $S \subset gH$. Similarly, for any $\theta \in H$ we have $g\circ \theta (B) = g (B) =A$, so $g\theta \in S$ or $gH \subset S$. Therefore, $gH =S$.

%\thang{Is it correct that in this case $S$ will be a full coset? or a subgroup? If yes, you should mention and include a proof.}

For small sets, we first look at an example, which says that $I(P, S)$ could be $|P||S|$. Let $S$ be a subset of matrices $\theta$ in $SL_2(\mathbb{F}_p)$ such that $\theta(0, 1)=(1, 0)$. So the size of $S$ can be arbitrary smaller than $p$. Let $P=\{ \lambda(e_1, e_2)\colon \lambda\in \mathbb{F}_p^*\}$, where $e_1=(1, 0), e_2=(0, 1)$. Then the size of $P$ can be arbitrary smaller than $p$ by choosing $\lambda$. With these sets $P$ and $S$, we have $I(P, S)=|P||S|$.

When we know better about structures of $B$ and $S$, then the following theorem is attained. 

\begin{theorem}\label{incidence-energy'}
For $\gamma\in (0, 1)$, there exists $\epsilon=\epsilon(\gamma)>0$ such that the following holds.

    Let $p$ be a sufficiently large prime. Let $P =A \times B \subseteq \l( \Fbb_p^2 \times \Fbb_p^2 \r) \setminus \lo (0, 0, 0, 0)\ro $, and let $S$ be a symmetric subset of $\SL_2 \l( \Fbb_p \r)$ such that $p^\gamma < \lv S\rv < p^{3-2\gamma} $ and $\lv S\cap gH \rv < p^{\frac{-\gamma}{2}}\lv S\rv$ for any subgroup $H \subsetneq \SL_2 \l( \Fbb_p \r)$ and $g \in \SL_2 \l( \Fbb_p \r)$. 
    \begin{enumerate}
        \item Assume $|B|\le p$, any line passing through the origin contains at most $k_1$ points from $B$, and $B$ determines at most $k_2$ distinct directions through the origin, then
        \[ I(P,S) \lesssim k_1^{\frac{1}{2}}|A|^{\frac{1}{2}}|S|+  \frac{k_1^{\frac{1}{2}}\lv B\rv^{\frac{1}{2}}|A|^{\frac{1}{2}}|S|^{\frac{3}{4}}}{p^{\frac{\epsilon}{4}}} + \frac{k_1^{\frac{1}{8}}\lv B\rv^{\frac{3}{4}}|A|^{\frac{1}{2}} |S|^{\frac{3}{4}}}{p^{\frac{\epsilon}{4}}}+\frac{k_1^{\frac{1}{4}}k_2^{\frac{1}{4}}\lv B\rv^{\frac{1}{2}} |A|^{\frac{1}{2}}|S|^{\frac{3}{4}}}{p^{\frac{\epsilon}{4}}}.\]
        \item Assume $|B|\le p^{8/15}$, and any line passing through the origin contains at most $k_1$ points from $B$, then
        \[ I(P,S) \lesssim  k_1^{\frac{1}{2}}|A|^{\frac{1}{2}}|S|+  \frac{k_1^{\frac{1}{2}}\lv B\rv^{\frac{1}{2}}|A|^{\frac{1}{2}}|S|^{\frac{3}{4}}}{p^{\frac{\epsilon}{4}}}+\frac{k_1^{\frac{1}{15}}|B|^{\frac{187}{225}} |A|^{\frac{1}{2}} |S|^{\frac{3}{4}}}{p^{\frac{\epsilon}{4}}}.\]
    \end{enumerate}
\end{theorem}
We note that the bound of $I(P, S)$ in this theorem is smaller than $|P||S|$ provided that $|A|\gg 1$.

\begin{remark}\label{rm24}
    If we remove the factor $p^{\epsilon/4}$ in the incidence bounds of Theorem \ref{incidence-energy'}, then the conditions on the set $S$ are not required.
\end{remark}

%\thang{We need to discuss about its sharpness}

%Compare with the trivial bound
%\[ I(P,S) \le |P||S|. \]
%For any $B$, we have
%\begin{align*}
%    k_1^{\frac{1}{2}}|A|^{\frac{1}{2}}|S| & \le |B||A||S| = |P||S|, \\
%    \frac{k_1^{\frac{1}{2}}|B|^{\frac{1}{2}}|A|^{\frac{1}{2}}|S|^{\frac{3}{4}}}{p^{\frac{\epsilon}{4}}} &\le |B||A|^{\frac{1}{2}}|S|^{\frac{3}{4}} \le |P||S|,\\
%    \frac{k_1^{\frac{1}{8}}\lv B\rv^{\frac{3}{4}}|A|^{\frac{1}{2}} |S|^{\frac{3}{4}}}{p^{\frac{\epsilon}{4}}} & \le |B||A|^{\frac{1}{2}}|S|^{\frac{3}{4}} \le |P||S|,\\
%    \frac{k_1^{\frac{1}{4}}k_2^{\frac{1}{4}}\lv B\rv^{\frac{1}{2}} |A|^{\frac{1}{2}}|S|^{\frac{3}{4}}}{p^{\frac{\epsilon}{4}}} & \le |B||A|^{\frac{1}{2}}|S|^{\frac{3}{4}} \le |P||S| \\
%    \frac{k_1^{\frac{1}{15}}|B|^{\frac{187}{225}} |A|^{\frac{1}{2}} |S|^{\frac{3}{4}}}{p^{\frac{\epsilon}{4}}} & \le |B||A|^{\frac{1}{2}}|S| \le |P||S|.
%\end{align*}
%So, the bound of Theorem \ref{incidence-energy'} is better than the trivial bound for all cases.

\subsection{Incidence mixing in the $SL_2(\mathbb{F}_p)$ action graph (Theorem \ref{incidence})}
% \label{section2}

Theorem \ref{incidence} will be proved by using tools from discrete Fourier analysis. We first recall some basic notations.

For $n\in \Nbb$, let $f:\Fbb_p^n \rar \Cbb$ be a complex valued function. The Fourier transform of $f$, denoted by $\widehat{f}$, is defined by
\[\widehat{f}(m):=p^{-n}\sum_{x\in\Fbb_p^n} \chi (-m\cdot x)f(x),\]
where $\chi$ is a nontrivial additive character of $\Fbb_p$. We
have the following basic properties of $\widehat{f}$.
\begin{itemize}
    \item The orthogonality property:
\[ \sum_{\alpha \in \Fbb_p^n}\chi (\beta \cdot \alpha) =\begin{cases}
    0 ,& \text{if}\quad \beta \ne (0,\ldots ,0) ,\\
    p^n ,& \text{if}\quad \beta = (0,\ldots ,0) .
\end{cases}\]
\item The Fourier inversion formula:
    \[ f(x)= \sum_{m\in \Fbb_p^n}\chi (m\cdot x) \widehat{f}(m). \]
\item The Plancherel formula:
    \[ \sum_{m\in \Fbb_p^n} \lv \widehat{f} (m)\rv^2 =p^{-n} \sum_{x\in \Fbb_p^n} \lv f(x) \rv^2. \]
    % \item Fourier transform of convolution of two funtion:
    % \[ \widehat{f \star g}(m) = \widehat{f}(m)\widehat{g}(m) \]
\end{itemize}

For $A\subseteq \Fbb_p^d $, by abuse of notation, we also denote its characteristic function by $A(x)$, i.e. $A(x)=1$ if $x\in A$ and $A(x)=0$ if $x\not\in A$.  

% \begin{theorem}[Theorem $2.6$, \cite{P.Y.23}]
%      Let $A,B\subseteq \Fbb_p^2 $ and $q\equiv 3 \mod 4$. We have
%     \[ \sum_{\lV m\rV = \lV m'\rV} \lv \widehat{A}(m) \rv^2 \lv \widehat{B}\l( m' \r) \rv^2 \lesssim \frac{\lv A\rv \lv B\rv}{p^{5}}\cdot \min \lo \lv A\rv^{\frac{1}{2}} ,\lv B\rv^{\frac{1}{2}} \ro. \]
% \end{theorem}

To proceed further, we need three lemmas. For each $x =\l( x_1,x_2 \r) \in \Fbb_p^2$, we define $x^\perp = \l( -x_2,x_1 \r)$. Note that $x\cdot y^\perp$ measures the area of the triangle with three vertices $x$, $y$, and the origin.  The first lemma presents a fact that the group $SL_2(\mathbb{F}_p)$ preserves areas of triangles with one vertex pinned at the origin.
\begin{lemma}\label{lmmay2}
Let $x, y, u, v$ be points in $\mathbb{F}_p^2\setminus\{(0, 0)\}$. If there exists $\theta \in \SL_2(\mathbb{F}_p)$ such that $\theta x=u$ and $\theta y=v$, then $x\cdot y^\perp=u\cdot v^\perp$. In the inverse direction, if $x$ and $y$ are not on the same line passing through the origin and $x\cdot y^\perp=u\cdot v^\perp$, then there exists unique $\theta\in SL_2(\mathbb{F}_p)$ such that $\theta x=u$ and $\theta y=v$. 
\end{lemma}
\begin{proof}
    Assume there exists $\theta \in \SL_2(\mathbb{F}_p)$ such that $\theta x=u$ and $\theta y=v$. Writing $\theta$ in the form 
    \[ \theta = \begin{bmatrix}
        a & b \\
        c & d
    \end{bmatrix} ,\]
    where $ad-bc=1$, $x=(x_1 ,x_2)$, and $y= (y_1 ,y_2)$. We have $u=(ax_1 +bx_2,cx_1 +dx_2) ,v = (ay_1 +by_2,cy_1 +dy_2) $. Then, 
    \begin{align*}
        u\cdot v^\perp &= -(ax_1+bx_2)(cy_1+dy_2) + (cx_1+dx_2)(ay_1+by_2) \\
        &= -(ad-bc)x_1y_2 +(ad-bc)x_2y_1 = x\cdot y^\perp .
    \end{align*}

In the inverse direction, since $x \ne ky$ for all $k \in \Fbb_p$, we have $x\cdot y^\perp \ne 0$. Indeed, writing $x$ as $x=\l( x_1,x_2 \r) \ne (0, 0)$. Since $x \ne (0, 0)$, we can assume that $x_1 \ne 0$. Therefore, if $y=\l( y_1,y_2 \r)$ satisfies $x\cdot y^\perp =0$, then $x$ and $y$ belong to a line passing through the origin, a contradiction. Similarly, we obtain $u\cdot v^\perp\ne 0$. Let $\theta$ be the matrix that maps the basis  $\lo x,y\ro$ to the basis $\lo u,v\ro$. We show that $\theta \in \SL_2(\mathbb{F}_p)$. Indeed, we write $\theta = \begin{bmatrix}
        a & b\\ 
        c & d
    \end{bmatrix}$ and $x= (x_1 ,x_2)$, and $y=(y_1 ,y_2)$. This implies $u=(ax_1 +bx_2,cx_1 +dx_2)$ and $v = (ay_1 +by_2,cy_1 +dy_2) $. Therefore, $x \cdot y^\perp = u\cdot v^\perp = (ad-bc)x\cdot y^\perp$, so $ad-bc=1$. In other words, $\theta\in SL_2(\mathbb{F}_p)$. 
\end{proof}

The next lemma tells us that the action of $SL_2(\mathbb{F}_p)$ on the plane $\mathbb{F}_p^2$ is transitive with multiplicity of $\sim p$. We give a general proof for the case $SL_n(\mathbb{F}_p)$. 
\begin{lemma}\label{lmmay1}
    For any $m,m'\in \Fbb_p^2 \setminus \lo (0,0)\ro $, define
    \[\M_{2,p}\l( m, m' \r) := \lo T\in \SL_2 \l( \Fbb_p \r) \colon Tm=m' \ro . \]
   Then  $|\M_{2,p}\l( m, m' \r)|= p.$
\end{lemma}

\begin{proof}
It follows from \cite[Theorem 13.3.3]{M.P.13} that $\SL_2(\mathbb{F}_p)$ is a group of size
    \[\lv \SL_2  \l( \Fbb_p \r) \rv = p\l( p^2-1 \r) .\]
Considering the group action of $\SL_2 \l( \Fbb_p \r) $ on $\Fbb_p^2$, $A : x \mapsto Ax , \forall A \in \SL_2 \l( \Fbb_p \r) ,x\in \Fbb_p^2$. For $m \ne (0,0)$, let $\text{Orb}(m)$ be the orbit of $m$, i.e. $\text{Orb}(m)=\{Tm\colon T\in \SL_2(\Fbb_p)\}$. Since $m\ne (0,0)$, there exists $m_1' \in \Fbb_p^2 \setminus \{ (0,0) \}$ such that 
$\lo m, m_1'\ro$ forms a linear independent system. So, let  
\[ T_m' = \begin{bmatrix}
    m&  m_1'
\end{bmatrix},\]
we have $\det T_m' = \lambda_m \ne 0$. Let $m_1= \lambda_m^{-1}m_1'$,  and
\[ T_m = \begin{bmatrix}
    m&  m_1
\end{bmatrix}.\]
Then,  $\det T_m = \lambda_m \cdot \lambda_m^{-1}=1$, so $T_m \in \SL_2(\Fbb_p)$. Now, for all $m' \in \Fbb_p^2\setminus \lo (0,0) \ro$, we observe that $T_{m'} \circ (T_m)^{-1}\in \SL_2(\Fbb_p)$, and $T_{m'} \circ (T_m)^{-1}(m)=m'$ so $\M_{2,p}(m,m') \ne \emptyset , \forall \, m,m' \ne (0,0)$. 

For $m, m'\in \mathbb{F}_p^2\setminus \{(0,0)\}$, let $T$ be an element in $\M_{2,p}(m, m')$. Then for all $T'\in \M_{2,p}(m, m')$, there exists $A\in \text{Stab}(m)$ such that $TA=T'$. This implies that $|\M_{2,p}\l( m,m' \r)|=\lv \text{Stab}(m) \rv $ for all $m,m'\in \Fbb_p^2 \setminus \{(0,0)\} $.

Moreover, for any $m\ne (0,0)$, there is no $T \in \SL_2(\Fbb_p)$ such that $T m =(0,0)$. Hence, we have $\text{Orb}(m) = \Fbb_p^2 \setminus \lo (0,0)\ro$. Thus, 
by the Orbit-Stabilizer theorem, 
\[ \lv \text{Stab}(m) \rv =\frac{\lv \SL_2 \l( \Fbb_p \r) \rv}{\lv \text{Orb} (m)\rv} = \frac{p(p^2-1)}{p^2-1}=p .\]This completes the proof. 
\end{proof}


The next lemma is an $L^2$ bound for the dot-product function.

\begin{lemma}[\cite{P.V.17}]\label{lmmay3}
    Let $A$ and $B$ be subsets of $\, \mathbb{F}_p^2$. The number of tuples $(x_1, x_2, y_1, y_2)\in A\times A\times B\times B$ such that $x_1\cdot x_2^{\perp}=y_1\cdot y_2^{\perp}$ is at most 
    \begin{equation}\label{eq111}\frac{|A|^2|B|^2}{p}+Cp^2|A||B|,\end{equation}
    for some positive constant $C$. 
\end{lemma}
% \begin{lemma}[\cite{P.V.17}]\label{lmmay3}
%     Let $A,B,C,D$ be subsets of $\, \mathbb{F}_p^2$. The number of tuples $(x_1, x_2, y_1, y_2)\in A\times B\times C\times D$ such that $x_1\cdot x_2=y_1\cdot y_2$ is at most 
%     \[\frac{|A||B||C||D|}{q}+Cp^2\sqrt{|A||B||C||D|},\]
%     for some positive constant $C$. 
% \end{lemma}
When $P\subset \mathbb{F}_p^4$ is a general set, the above upper bound can be replaced by $\frac{|P|^2}{p}+Cp^2|P|$. Let $k_A$ and  $k_B$ be the maximal number of points from $A$ and $B$ on a line passing through the origin, respectively. Assume that $\min\{k_A, k_B\}=k$, then, with the same argument, the bound (\ref{eq111}) can be replaced by
    \[\frac{|A|^2|B|^2}{p}+Cpk|A||B|.\]

% \begin{lemma}[\cite{P.V.17}]\label{lmmay9}
%     Let $A$ and $B$ be subsets of $\, \mathbb{F}_p^2$. The number of tuples $(x_1, x_2, y_1, y_2)\in A\times A\times B\times B$ such that $x_1\cdot x_2^perp =y_1\cdot y_2^perp$ is at most 
%     \[\frac{|A|^2|B|^2}{q}+Cp^2|A||B|,\]
%     for some positive constant $C$. When $P\subset \mathbb{F}_p^4$ is a general set, the above upper bound can be replaced by $|P|^2/q+Cp^2|P|$. If we assume in addition that any line passing through the origin contains at most $k$ points from $A$, then the number of the tuples is at most 
%     \[\frac{|A|^2|B|^2}{q}+Cqk|A||B|.\]
% \end{lemma}

With these three lemmas in hand, we are ready to prove Theorem \ref{incidence}.

\begin{proof}[Proof of Theorem \ref{incidence}]
Without loss of generality, we assume that $(0, 0, 0, 0)\not\in P$, since this element only contributes $|S|$ incidences to the incidence bound. 

    By using the Fourier transformation and the Fourier inversion formula, we have 
    \begin{align*}
        I(P,S) & = \sum_{\substack{p=(x,y )\in P \\ \theta \in S}} 1_{\theta y=x} =\frac{1}{p^2} \sum_{m\in \Fbb_p^2} \sum_{\substack{(x,y)\in P, \\ \theta \in S}} \chi \l( m\cdot \l( x-\theta y \r) \r) \\
        & = \frac{\lv P\rv \lv S\rv}{p^2} + \frac{1}{p^2} \sum_{ m\in \Fbb_p^2 \setminus \lo (0,0)\ro } \sum_{\substack{ (x,y)\in P ,\\ \theta \in S }} \chi \l( m \cdot \l( x -\theta y \r) \r) \\
        & = \frac{\lv P\rv \lv S\rv}{p^2} + p^2\sum_{m\ne (0, 0)} \sum_{\theta \in S} \widehat{P} \l( -m,\theta^t m\r).
    \end{align*}
By using the Cauchy--Schwarz inequality, one has 
\begin{align*}
    &\sum_{\theta\in S}\sum_{m\ne (0, 0)}\widehat{P}(-m, \theta^Tm)\le |S|^{\frac{1}{2}}\left(\sum_{\theta\in \SL_2 \l( \Fbb_p \r)}\sum_{m_1, m_2\ne (0, 0)}\widehat{P}(-m_1, \theta^tm_1)\overline{\widehat{P}(-m_2, \theta^tm_2)}\right)^{\frac{1}{2}}.
\end{align*}
We now observe 
\begin{align*}
    &\sum_{\theta\in \SL_2 \l( \Fbb_p \r)}\sum_{m_1, m_2\ne (0, 0)}\widehat{P}(-m_1, \theta^tm_1)\overline{\widehat{P}(-m_2, \theta^tm_2)}\\&=\frac{1}{p^8}\sum_{\theta}\sum_{m_1, m_2\ne (0, 0)}\sum_{(x_1, y_1), (x_2, y_2)}P(x_1, y_1)P(x_2, y_2)\chi(-m_1x_1+\theta^tm_1y_1)\chi(m_2x_2-\theta^tm_2y_2)\\
    &=\sum_{m_1, m_2}-\sum_{m_1=(0, 0), m_2\ne (0, 0)}-\sum_{m_1\ne (0, 0), m_2=(0, 0)}-\sum_{m_1=m_2=(0, 0)}\\
    &=:I-II-III-IV.
\end{align*}
We now estimate each term separately. 
\begin{align*}
    &I=\frac{1}{p^8}\sum_{\theta}\sum_{m_1, m_2\in \mathbb{F}_p^2}\sum_{(x_1, y_1), (x_2, y_2)}P(x_1, y_1)P(x_2, y_2)\chi(m_1(\theta y_1-x_1))\chi(m_2(\theta y_2-x_2))\\
    &=\frac{1}{p^4}\sum_{\theta}\sum_{(x_1, y_1), (x_2, y_2)}P(x_1, y_1)P(x_2, y_2)1_{\theta y_1=x_1}1_{\theta y_2=x_2}.\\
\end{align*}
By using  Lemma \ref{lmmay2}, one has
\begin{align*}
    &I=\frac{1}{p^4}\sum_{\theta}\sum_{\lambda\in \mathbb{F}_p}\sum_{\substack{x_1, y_1, x_2, y_2\\~y_1=\lambda y_2, x_1=\lambda x_2}}A(x_1)A(x_2)B(y_1)B(y_2)1_{\theta y_1=x_1}1_{\theta y_2=x_2}\\
    &+\frac{1}{p^4}\sum_{x_1, y_1, x_2, y_2}A(x_1)A(x_2)B(y_1)B(y_2)1_{y_1\cdot y_2^\perp=x_1\cdot x_2^\perp}.
\end{align*}
Notice that Lemma \ref{lmmay1} and Lemma \ref{lmmay2} tell us that the first sum can be bound by at most $p^2|A||B|/p^4$.

By Lemma \ref{lmmay3}, the second sum can be at most 
\[\frac{|A|^2|B|^2}{p^5}+C\frac{|A||B|}{p^2}.\]

In other words, we have 
\[I\le \frac{|A|^2|B|^2}{p^5}+(C+1)\frac{|A||B|}{p^2}.\]

Moreover,
\[IV=\frac{(p^3-p)|A|^2|B|^2}{p^8}=\frac{|A|^2|B|^2}{p^5}-\frac{|A|^2|B|^2}{p^7}.\]

Thus, $I-IV\ll \frac{|A||B|}{p^2}$.
Regarding $II$, 
\begin{align*}
    II & =\frac{1}{p^8}\sum_\theta \sum_{( x_1,y_1)} P(x_1,y_1) \l( \sum_{(x_2,y_2)} P(x_2,y_2) \l( \sum_{m_2\ne 0} \chi \l( m_2 \l( x_2 -\theta y_2 \r) \r) \r) \r) \\
    & = \frac{1}{p^8} \sum_\theta \sum_{(x_1,y_1)} P(x_1 ,y_1) \l( \sum_{\substack{(x_2,y_2), \\ x_2 =\theta y_2}} P(x_2,y_2) \l( p^2-1 \r)  - \sum_{ \substack{(x_2,y_2), \\ x_2 \ne \theta y_2 }} P(x_2,y_2)  \r) \\
    & = \frac{1}{p^8} \lv P\rv \l( \l( p^2-1\r)\sum_{\substack{(x_2,y_2) \in P,\\ \theta \in \SL_2 \l( \Fbb_p \r) ,\\ x_2 =\theta y_2}} 1  - \sum_{ \substack{(x_2,y_2) \in P,\\ \theta \in \SL_2 \l( \Fbb_p \r) ,\\ x_2 \ne \theta y_2 }} 1\r) \\
    & = \frac{1}{p^8} \lv P\rv \l( \lv A\rv \lv B\rv p\l( p^2 -1\r) - \lv A\rv \lv B\rv \l( p^3 -2p \r) \r) = \frac{\lv P\rv^2}{p^7}.
\end{align*}

Similarly, we obtain  $III =\frac{\lv P\rv^2}{p^7} .$
Putting all estimates together, we conclude that 
\[\left\vert I(P, S)-\frac{|P||S|}{p^2}\right\vert\ll p\sqrt{|P||S|}+|S|.\]
This completes the proof.
\end{proof}

\subsection{Energy refinement and subgroup non-concentration (Theorem \ref{thm: 4.12})}\label{sectionenergy}
% \label{section2}

%\thang{This section is not complete yet!}
%One might wonder if the proofs of Theorem \ref{incidence} and Theorem \ref{incidence-energy} can be combined in a certain way to give a new type incidence estimate. In this section, we discuss such a question. 

For a set $S\subset SL_2(\mathbb{F}_p)$, we define the energy $\Ebf(S, S)$ by
\[\Ebf(S, S):=\#\{(a, b, c, d)\in S^4\colon ab=cd\}.\]
The trivial bound of $\Ebf(S, S)$ is $|S|^3$. When $S$ is a large set, Babai, Nikolov, and Pyber proved in \cite{B.N.P.08} that
 \begin{equation}\label{eq-energy1} \Ebf \l( S, S \r) \ll p^2 |S|^2 + \frac{|S|^4}{p^3} .\end{equation}
This bound is sharp. To see its sharpness, we provide an example here.

  Let $S$ be the set of matrices of the form
    \[ \begin{bmatrix}
        \ast & -x \\
        x^{-1} & 0
    \end{bmatrix}, \]
    where $x,\ast \in \Fbb_p \setminus \{ 0 \}$. Then, $S$ is a subset of $\SL_2 \l( \Fbb_p \r)$ and $\lv S\rv =(p-1)^2$. We consider following equation
    \begin{align}\label{eq: 3}
        \begin{bmatrix}
        \ast_1 & -x \\
        x^{-1} & 0
    \end{bmatrix} \begin{bmatrix}
        \ast_2 & -y \\
        y^{-1} & 0
    \end{bmatrix} =  \begin{bmatrix}
        \ast_1' & -x' \\
        (x')^{-1} & 0
    \end{bmatrix} 
    \begin{bmatrix}
        \ast_2' & -y' \\
        (y')^{-1} & 0
    \end{bmatrix},
    \end{align} 
    where all matrices are in $S$. The equation is equivalent to 
    \[ \begin{cases}
        \ast_1 \ast_2 -xy^{-1} = \ast_1' \ast_2 ' -x'(y')^{-1}, \\
        -y\ast_1 = -y' \ast_1' ,\\
        \ast_2 x^{-1} =\ast_2'(x')^{-1} ,\\
        -yx^{-1} =-y' x^{-1}.
    \end{cases} \]
    This implies $\frac{y}{y'}=\frac{x}{x'}=\frac{\ast_2}{\ast_2'}=\frac{\ast_1'}{\ast_1}$. Therefore, for fixed $\ast_1 , \ast_1' ,x,x'$, there exist $(p-1)^2$ tuples $\l( y,y',\ast_2,\ast_2' \r) \in \l( \Fbb_p \setminus \{ 0\} \r)^4$ such that the equation (\ref{eq: 3}) holds. In other words, for each pair of matrices $(A,C) \in S^2$, there exist $(p-1)^2$ pairs of matrices $(B,D) \in S^2$ such that $AB=CD$. Hence,
    \[ \Ebf (S,S) \gg \lv S\rv^2 (p-1)^2 =(p-1)^6 \sim p^2 (p-1)^4 + \frac{(p-1)^8}{p^3} = p^2\lv S\rv^2 + \frac{\lv S\rv^4}{p^3}. \]
When the set $S$ is of small size, one would hope to have an upper bound of $\Ebf(S, S)$ that does not depend on $p$. In this paper, we make use of the following result due to Bourgain and Gamburd in \cite{B.A.08} to derive such a bound over prime fields.
\begin{theorem}[Proposition $2$, \cite{B.A.08}]\label{thm: 7.1}
Let $p$ be a sufficiently large prime, and let $\eta$ be a symmetric probability measure on $\SL_2(\Fbb_p)$ and $0< \gamma < \frac{3}{4}$, such that
\begin{itemize}
\item[(1)] $\lV \eta  \rV_{\infty} < p^{-\gamma} ;$
\item[(2)] $\eta(gH)<p^{\frac{-\gamma }{2}} $ for any proper subgroup $H\subset SL_2( \Fbb_p ), g\in SL_2( \Fbb_p );$
\item[(3)] $\Vert \eta\Vert_2> p^{\frac{-3}{2} +\gamma}$.
\end{itemize}
Then there exists $\epsilon = \epsilon \l( \gamma \r) >0$ such that
$$
\Vert\eta*\eta\Vert_2 < p^{-\epsilon} \Vert\eta\Vert_2.
$$
\end{theorem}

Moreover, following the proof of Theorem \ref{thm: 7.1}, we have $\epsilon < \gamma$. 

For $0< \gamma <\frac{3}{4}$, let $S$ be a symmetric subset of $\SL_2 \l( \Fbb_p \r)$ such that $p^{\gamma}<|S|<p^{3-2\gamma}$ and for any proper subgroup $H \subsetneq \SL_2 \l( \Fbb_p \r) $, $g \in \SL_2 \l( \Fbb_p \r)$ we have $\lv S\cap gH\rv < p^{\frac{-\gamma }{2}}\lv S\rv $. Let $\mu_S : \SL_2 \l( \Fbb_p \r) \rar \Rbb $ be the function defined by $\mu_S \l( g \r) =\frac{1}{\lv S\rv}$ if $g\in S$ and $\mu (g)=0$ if $g \notin S$. Then, $\mu_S$ satisfies all conditions of Theorem \ref{thm: 7.1}. Indeed, we have 
\begin{itemize}
    \item[(1)] $\lV \mu_S \rV_\infty = \max_{g} \mu_S(g) = \frac{1}{|S|} < p^{-\gamma }$,
    \item[(2)] $\mu_S (gH) = \frac{|S\cap gH|}{|S|} < p^{\frac{-\gamma }{2}}$, for any proper subgroup $H \subset \SL_2(\Fbb_p)$ and $g\in \SL_2(\Fbb_p)$,
    \item[(3)] $\lV \mu_S \rV_2 = \l( \sum_{g\in \SL_2(\Fbb_p)} \mu_S(g)^2 \r)^{\frac{1}{2}} = \frac{1}{|S|^{\frac{1}{2}}} > p^{\frac{-3}{2}+\gamma}.$
\end{itemize}
Therefore,
\[ \lV \mu_S \ast \mu_S \rV_2 < p^{-\epsilon} \lV \mu_S \rV_2 = p^{-\epsilon} \frac{1}{\lv S\rv^{\frac{1}{2}}}. \]

One the other hand, $\mu_S \ast \mu_S (g) = \frac{\lv  \lo (a,b)\in S \times S \colon ab=g \ro \rv}{\lv S\rv^2} $. Then,
\[ \lV \mu_S \ast \mu_S \rV_2^2 = \frac{\Ebf \l( S,S \r)}{\lv S\rv^4}.\]
In other words, we have proved the following corollary. 
\begin{corollary}\label{co:36}
For any $\gamma\in (0, 1)$, there exists $\epsilon=\epsilon(\gamma)>0$ such that the following holds.
 Let $p$ be a sufficiently large prime, let $S$ be a symmetric subset of $\SL_2 (\Fbb_p )$ such that $p^\gamma < |S| < p^{3-2\gamma } $ and $ |S \cap gH| < p^{\frac{-\gamma}{2}}|S|$ for any subgroup $H \subsetneq \SL_2 (\Fbb_p) $ and $g\in \SL_2(\Fbb_p)$. Then, we  have
\begin{equation}
\Ebf \l( S,S \r) < \frac{\lv S\rv^3}{p^{\epsilon}}.
\end{equation}    
\end{corollary}
\begin{proof}
    For any $\gamma\in (0, \frac{3}{4})$, it follows directly from the previous computation. For any $\gamma\in [3/4, 1)$, we observe that 
    \[(p^\gamma, p^{3-2\gamma})\subset (p^{\gamma'}, p^{3-2\gamma'}), ~~p^{-\frac{\gamma}{2}}|S|\le p^{-\frac{\gamma'}{2}}|S|,\]
    for any $\gamma'\in (0, 3/4)$. Then, by choosing $\epsilon=\epsilon(\gamma')$ for any $\gamma'\in (0, 3/4)$, the corollary follows.
\end{proof}


 Let $A \times B \subseteq \Fbb_p^2 \times \Fbb_p^2 $ and $S\subseteq \SL_2 \l( \Fbb_p \r)$. Then,
    \[ I(A\times B,S) \le N^{\frac{1}{2}}\lv A\rv^{\frac{1}{2}} ,\]
    where $N$ is the number of $\l( b,b',\theta ,\theta' \r) \in B \times B \times S \times S$ such that $\theta b  =\theta 'b'$.

Indeed,  for each $\l( a,b \r) \in A\times B$, denote $s_{(a,b)} $ as the number of $\theta \in S$ such that $\theta b=a$. Therefore,
    \begin{align*}
        I(P,S) & = \sum_{a\in A} \l( \sum_{b\in B} s_{(a,b)} \r)  \le  \lv A\rv^{\frac{1}{2}} \l( \sum_{a\in A} \l( \sum_{b\in B} s_{(a,b)} \r)^2 \r)^{\frac{1}{2}} \\
        & =  \lv A\rv^\frac{1}{2} \l( \sum_{a\in A} \lv \lo \l( b,b',\theta ,\theta' \r) \in B\times B \times S \times S \colon \theta b=\theta 'b' =a \ro \rv \r)^{\frac{1}{2}} \\
        & \le \lv A\rv^{\frac{1}{2}} \cdot N^{\frac{1}{2}}.
    \end{align*}
To bound $N$, we observe that the equation $\theta b=\theta'b'$ gives $(\theta')^{-1}\theta b=b'$. So, $N$ can be viewed as the number of incidences between $B\times B$ and the multi-set $S^{-1}S$.  If $(0, 0) \notin B$, by following the proof of Theorem \ref{incidence} identically, we obtain
\[N\ll  \frac{\lv B\rv^2 \lv S\rv^2}{p^2} +p\lv B\rv (\Ebf (S,S))^{\frac{1}{2}},\]
and if any line passing through the origin contains at most $k$ points from $B$, then
\[ N \ll  \frac{\lv B\rv^2 \lv S\rv^2}{p^2} +p^{\frac{1}{2}}k^{\frac{1}{2}}\lv B\rv (\Ebf (S,S))^{\frac{1}{2}}.\]

As mentioned in (\ref{eq-energy1}), 
 \[ \Ebf \l( S, S \r) \ll p^2 |S|^2 + \frac{|S|^4}{p^3} .\]
Substituting this bound into $I(A\times B, S)$ implies 
 \[ I(A\times B ,S)\ll \frac{|A|^{\frac{1}{2}}|B||S|}{p}+\frac{|A|^{\frac{1}{2}}|B|^{\frac{1}{2}}|S|}{p^{\frac{1}{4}}}+p|A|^{\frac{1}{2}}|S|^{\frac{1}{2}}|B|^{\frac{1}{2}}.\]

Compared to the bound of Theorem \ref{incidence}, this result is weaker. 

However, if we use Corollary \ref{co:36} instead, then
    \[ I(P,S) \ll \frac{|A|^{\frac{1}{2}}|B||S|}{p} + p^{\frac{2-\epsilon}{4}}|A|^{\frac{1}{2}}|B|^{\frac{1}{2}}|S|^{\frac{3}{4}}.\]
Moreover,  if any line passing through the origin contains at most $k$ points from $B$, then
    \[ I(P,S) \ll \frac{|A|^{\frac{1}{2}}|B||S|}{p} + k^{\frac{1}{4}}p^{\frac{1-\epsilon}{4}}|A|^{\frac{1}{2}}|B|^{\frac{1}{2}}|S|^{\frac{3}{4}}.\]
This completes the proof of Theorem \ref{thm: 4.12}.



\subsection{Two alternative approaches yield weaker bounds}
% \label{section2}
This section presents two different approaches without techniques from Fourier analysis. Although the resulting bounds are weaker compared to those of Theorem \ref{incidence} and Theorem \ref{thm: 4.12}, the methods will be useful for us when studying the case of small sets. 

%\thang{In this section, we mention two other approaches, but weaker bounds for large sets}
\begin{theorem}\label{thm:3.7}
For any $\gamma\in (0, 1)$, there exists $\epsilon=\epsilon(\gamma)>0$ such that the following holds.
    Let $p$ be a sufficiently large prime. Let $P=A\times B \subseteq (\Fbb_p^2 \times \Fbb_p^2 ) \setminus \lo (0, 0, 0, 0)\ro$, and let $S$ be a symmetric subset of $\SL_2 (\Fbb_p )$ such that $p^\gamma < |S| < p^{3-2\gamma } $ and $ |S \cap gH| < p^{\frac{-\gamma }{2}}|S|$ for any subgroup $H \subsetneq \SL_2 (\Fbb_p) $ and $g\in \SL_2(\Fbb_p)$. Assume any line passing through the origin contains at most $k$ points from $B$. Then, 
    \begin{itemize}
        \item[(1)] if $|B|< k^{\frac{1}{2}}p$, we have 
        \[ I(P,S) \lesssim k^{\frac{1}{2}}|A|^{\frac{1}{2}} |S| + k^{\frac{1}{4}} p^{\frac{1-\epsilon}{4}} |A|^{\frac{1}{2}} |B|^{\frac{1}{2}} |S|^{\frac{3}{4}} ,\]
        \item[(2)] if $|B|\ge k^{\frac{1}{2}}p$, we have 
        \[ I(P,S) \lesssim 
        % k^{\frac{1}{2}}|A|^{\frac{1}{2}} |S| +
        \frac{|A|^{\frac{1}{2}}|B||S|^{\frac{3}{4}}}{p^{\frac{1+ \epsilon }{4}}}.\]
    \end{itemize}
    % \[ I(P,S) \lesssim k^{\frac{1}{2}}|A|^{\frac{1}{2}} |S| +  p^{\frac{-1-\epsilon}{4}}|A|^{\frac{1}{2}}|B||S|^{\frac{3}{4}} + k^{\frac{1}{4}} p^{\frac{1-\epsilon}{4}} |A|^{\frac{1}{2}} |B|^{\frac{1}{2}} |S|^{\frac{3}{4}} ,\]
   
\end{theorem}
%\thang{Hung, can you add a comparison between this theorem and Theorem \ref{incidence} and Theorem \ref{thm: 4.12}?}
Compared to the
% the bound in Theorem \ref{thm:3.7} and the 
bound of Theorem \ref{thm: 4.12}, which is
\[ I(P,S) \ll \frac{|A|^{\frac{1}{2}}|B||S|}{p} + k^{\frac{1}{4}}p^{\frac{1-\epsilon}{4}}|A|^{\frac{1}{2}}|B|^{\frac{1}{2}}|S|^{\frac{3}{4}},\]
we can see that
\begin{itemize}
    \item if $|B|< k^{\frac{1}{2}}p $, then
    \begin{align*}
    \frac{|A|^\frac{1}{2}|B||S|}{p}< k^{\frac{1}{2}}|A|^{\frac{1}{2}}|S|,
\end{align*}
    \item if $|B|\ge k^{\frac{1}{2}}p $, then
    \begin{align*}
    \frac{|A|^\frac{1}{2}|B||S|}{p}, k^{\frac{1}{4}}p^{\frac{1-\epsilon}{4}}|A|^{\frac{1}{2}}|B|^\frac{1}{2}|S|^{\frac{3}{4}} \le \frac{|A|^{\frac{1}{2}}|B||S|^{\frac{3}{4}} }{p^{\frac{1+\epsilon}{4}}},
\end{align*}
since $|S|< p^{3-2\gamma} < p^{3-\epsilon}$. In other words, Theorem \ref{thm: 4.12} is better than Theorem \ref{thm:3.7}.
\end{itemize}


\begin{proof}[Proof of Theorem \ref{thm:3.7}]
Since the identity matrix $I$ contributes at most $$\min \lo |A|,|B|\ro \ll k^{\frac{1}{4}}p^{\frac{1-\epsilon }{4}}\lv A\rv^{\frac{1}{2}} \lv B\rv^{\frac{1}{2}} |S|^{\frac{3}{4}}$$ incidences to $I(P,S)$, we may assume without loss generality that $I \not\in S$.

    For a set $D\subset \mathbb{F}_p^2\times \mathbb{F}_p^2$ and $\theta\in SL_2(\mathbb{F}_p)$, by $i_{D}(\theta)$, we mean the number of incidences between $\theta$ and the set $D$. 

We have 
\begin{align*}
I(P, S)=\sum_{\theta\in S}i_P(\theta) \le |A|^{\frac{1}{2}}\left(I(B\times B, S')\right)^{\frac{1}{2}},
\end{align*}
where $S'$ be the multi-set of elements of the form $a^{-1}b$, where $a,b\in S$. Set $E=B\times B$. For $\theta\in S'$, let $m(\theta)$ be the multiplicity of $\theta$ in $S'$, and for $(x, y)\in E$, let $n(x, y)$ be the number of elements $(u, v)\in E$ such that $(u, v)=\lambda (x, y)$ for some $\lambda\in \mathbb{F}^*_p$. 

We now bound the number of incidences between $S'$ and $B\times B$. By $\overline{S'}$, we mean the set of distinct elements in $S$. We have 
\[I(B\times B, S')=\sum_{\theta\in \overline{S'}}m(\theta)\sum_{(x, y)\in E}\theta(x, y),\]
where $\theta(x, y)=1$ if $\theta y=x$, and $0$ otherwise. 

We now write
\begin{align*}
    &I(B\times B, S')=\sum_{\theta\in \overline{S'}}m(\theta)\sum_{(x, y)\in E}\theta(x, y)\\
    &=\sum_{\theta\in \overline{S'}}m(\theta)\sum_{(x, y)\in E, n(x, y)=1}\theta(x, y)+\sum_{\theta\in \overline{S'}}m(\theta)\sum_{(x, y)\in E, n(x, y)>1}\theta(x, y)=I+II.
\end{align*}

Let $E'$ be the set of $(x, y)\in E$ such that $n(x, y)=1$. 
To bound $I$, 

\begin{align*}
    I=\sum_{\theta\in \overline{S'}, ~i_{E'}(\theta)>1}m(\theta)\sum_{(x, y)\in E, n(x, y)=1}\theta(x, y)+\sum_{\theta\in \overline{S'}, ~i_{E'}(\theta)=1}m(\theta)\sum_{(x, y)\in E, n(x, y)=1}\theta(x, y)=I_1+I_2.
\end{align*}
It is clear that $I_2\le |S'|=|S|^2$. By the Cauchy--Schwarz inequality and Lemma \ref{lmmay3}, we have 
\begin{align*}
     I_1 & =\sum_{\theta \in\overline{S'},i_{E'}(\theta)>1}m (\theta )i_{E'}(\theta) \le  \l(\sum_{\theta \in\overline{S'},i_{E'}(\theta)>1}m (\theta )^2 \r)^\frac{1}{2} \l( \sum_{\theta \in\overline{S'},i_{E'}(\theta)>1} i_{E'}(\theta)^2 \r)^{\frac{1}{2}}   \\
     & \ll \l( \Ebf (S,S) \r)^{\frac{1}{2}} \l( \sum_{\theta \in\overline{S'},i_{E'}>1}\binom{i_{E'}(\theta)}{2} \r)^{\frac{1}{2}} \\
     & \le \l( \Ebf (S,S) \r)^{\frac{1}{2}} \l( \lv \lo  (x_1,y_1,x_2,y_2) \in E \colon (x_1,y_1)\ne (x_2,y_2), \exists \theta \in \overline{S'}, \theta \text{ is incident to } (x_1,y_1) \text{ and } (x_2,y_2) \ro \rv \r)^{\frac{1}{2}} \\
     & \le \l( \Ebf (S,S) \r)^{\frac{1}{2}} \l( \lo \l( x_1,y_1,x_2,y_2 \r) \in B^4 \colon \exists \theta \in \overline{S'} , \theta x_1 =y_1 ,\theta x_2 =y_2 \ro \r)^{\frac{1}{2}} \\
     & \le \l( \Ebf (S,S) \r)^{\frac{1}{2}} \l( \lo \l( x_1,y_1,x_2,y_2 \r) \in B^4 \colon x_1 \cdot x_2^\perp =y_1 \cdot y_2^\perp \ro \r)^{\frac{1}{2}} \\
     &\ll (\Ebf (S,S))^{\frac{1}{2}}\cdot \l( \frac{|B|^2}{p^{\frac{1}{2}}} + k^\frac{1}{2}p^\frac{1}{2}|B| \r) .
    \end{align*}
    
Thus, using Corollary \ref{co:36}, we obtain 
\[I\le |S|^2+ \l( \frac{|B|^2}{p^{\frac{1}{2}}} + k^{\frac{1}{2}}p^{\frac{1}{2}}|B|\r) |S|^{\frac{3}{2}}p^{\frac{-\epsilon}{2}}.\]
We now consider $II$. 
% \begin{align*}
%   II= \sum_{\theta\in \overline{S'}}m(\theta)\sum_{(x, y)\in E, n(x, y)>1}\theta(x, y)\sim \sum_{\theta\in \overline{S'}}m(\theta)\sum_{(x, y)\in E, n(x, y)\sim \ell}\theta(x, y).
% \end{align*}
\begin{align*}
    II &= \sum_{\theta\in \overline{S'}}m(\theta)\sum_{(x, y)\in E, n(x, y)>1}\theta(x, y)=\sum_{i=1}^{\lf \log_2 (k) \rf} \sum_{\theta\in \overline{S'}}m(\theta)\sum_{(x, y)\in E, 2^i \le n(x, y) < 2^{i+1}}\theta(x, y).
\end{align*}
By using pigeonhole principle, there exists $\ell =2^{i_0} $ such that
\[ II \lesssim  \sum_{\theta\in \overline{S'}}m(\theta)\sum_{(x, y)\in E, \ell \le n(x, y)< 2\ell}\theta(x, y) . \]
% \[ II =\Omega_{\log_2 (k)} \l( \sum_{\theta\in \overline{S'}}m(\theta)\sum_{(x, y)\in E, \ell \le n(x, y)< 2\ell}\theta(x, y) \r). \]
We say two pairs $(x, y)$ and $(x', y')$ are in the same congruence class if there exists $\lambda \in \Fbb_p^\ast$ such that $(x,y) = \lambda (x',y')$. Define $E''$ to be the set of congruence classes $[(x, y)]$ in $E$ such that $\ell \le n(x, y)< 2\ell$. 

Then we have 
\[II \lesssim \ell\sum_{\theta\in \overline{S'}}m(\theta)i_{E''}(\theta)=II_1+II_2.\]
% \[II\sim_{\log_2 (k)} \ell\sum_{\theta\in \overline{S'}}m(\theta)i_{E''}(\theta)=II_1+II_2.\]
Here 
\[II_1=\ell\sum_{\theta\in \overline{S'}, ~i_{E''}(\theta)=1}m(\theta)i_{E''}(\theta),~~~II_2=\ell\sum_{\theta\in \overline{S'}, ~i_{E''}(\theta)\ge 2}m(\theta)i_{E''}(\theta).\]
As above, we have 
\[II_1\le \ell |S|^2,\]
and 
\begin{align*}
    II_2 & \le \l( (\Ebf (S,S)) \r)^{\frac{1}{2}} \ell \l( \sum_{\theta\in \overline{S'}, ~i_{E''}(\theta)\ge 2} i_{E''}(\theta)^2 \r) ^{\frac{1}{2}}  \ll \l( \Ebf (S,S) \r)^\frac{1}{2} \l( \sum_{\theta \in\overline{S'},i_{E''}(\theta )_>1} \ell^2 \binom{i_{E''}(\theta)}{2} \r)^{\frac{1}{2}} \\
    & \le \l( \Ebf (S,S) \r)^\frac{1}{2} \l( \ell^2 \lv \lo  (x_1,y_1,x_2,y_2) \in E''\times E''\colon \exists \theta \in \overline{S'} , \theta \text{ is incident to $(x_1,y_1)$ and $(x_2,y_2)$} \ro \rv \r)^{\frac{1}{2}} \\
    & \le \l( \Ebf (S,S) \r)^\frac{1}{2} \l( \lv \lo  (x_1,y_1,x_2,y_2) \in B\times B \times B \times B\colon \exists \theta \in \overline{S'} , \theta y_1=x_1 ,\theta y_2 =x_2 \ro \rv \r)^{\frac{1}{2}} \\
    & \le \l( \Ebf (S,S) \r)^\frac{1}{2} \l( \lv \lo \l( x_1,y_1,x_2,y_2\r) \in B\times B \times B \times B \colon x_1\cdot x_2^\perp =y_1 \cdot y_2^\perp \ro \rv \r)^{\frac{1}{2}} \\
    & \le (\Ebf (S,S))^{\frac{1}{2}} \cdot \l( \frac{|B|^2}{p^{\frac{1}{2}}} + k^{\frac{1}{2}}p^{\frac{1}{2}}|B| \r) .
\end{align*} 
% the forth inequality follows from that for each $(x,y)\in (E'')^2$, there are $\ge \ell^2$ quadruples $(x_1,y_1,x_2,y_2) \in B^4$ such that $(x_1,y_1) = \lambda_1 x, (x_2,y_2) =\lambda_2 y$, for some $ \lambda_1 ,\lambda_2 \in \Fbb_p^\ast$.


Putting these bounds together implies 
\[II\lesssim \ell |S|^2+ (\Ebf (S,S))^{\frac{1}{2}} \l( \frac{|B|^2}{p^{\frac{1}{2}}} + k^{\frac{1}{2}}p^{\frac{1}{2}}|B| \r).\]
Notice that $\ell\le k$. So,
\[I(B\times B, S')\lesssim  k|S|^2+\l( \frac{|B|^2}{p^{\frac{1}{2}}} + k^{\frac{1}{2}}p^{\frac{1}{2}}|B| \r) |S|^{\frac{3}{2}}p^{\frac{-\epsilon}{2}},\]
and then
\[ I(P,S) \lesssim k^{\frac{1}{2}}|A|^{\frac{1}{2}} |S| + k^{\frac{1}{4}} p^{\frac{1-\epsilon}{4}} |A|^{\frac{1}{2}} |B|^{\frac{1}{2}} |S|^{\frac{3}{4}} + p^{\frac{-1-\epsilon}{4}}|A|^{\frac{1}{2}}|B||S|^{\frac{3}{4}} .\]

A direct computation implies that
\begin{enumerate}
    \item if $|B|< k^{\frac{1}{2}}p$, then
    \[ I(P,S) \lesssim k^{\frac{1}{2}}|A|^{\frac{1}{2}} |S| + k^{\frac{1}{4}} p^{\frac{1-\epsilon}{4}} |A|^{\frac{1}{2}} |B|^{\frac{1}{2}} |S|^{\frac{3}{4}} ;\]
    \item if $|B| \ge k^{\frac{1}{2}}p$, then
    \[ I(P,S) \lesssim  p^{\frac{-1-\epsilon}{4}}|A|^{\frac{1}{2}}|B||S|^{\frac{3}{4}}.\]
\end{enumerate}
This completes the proof.
\end{proof}


\begin{theorem}\label{thm: 3.8}
    Let $p$ be a prime and $P=A\times B\subseteq (\Fbb_p^2 \times \Fbb_p^2) \setminus \lo (0, 0, 0, 0)\ro $. Assume any line passing through the origin contains at most $k$ points from $B$. Then, we have
    \[ I(P,S) \lesssim   \frac{|A||B|\lv S\rv^{\frac{1}{2}}}{p^{\frac{1}{2}}} + k^{\frac{1}{2}}p^{\frac{1}{2}}|A|^{\frac{1}{2}}|B|^{\frac{1}{2}}\lv S\rv^{\frac{1}{2}} + k\lv S\rv .\]
    
\end{theorem}
Compared to the bound of Theorem \ref{incidence}, which is
\[ I(P,S) \ll \frac{|A||B||S|}{p^2} + k^{\frac{1}{2}}p^{\frac{1}{2}}|A|^{\frac{1}{2}}|B|^{\frac{1}{2}}|S|^{\frac{1}{2}},\]
we can see that
\[ \frac{|A||B||S|}{p^2} \le \frac{|A||B||S|^{\frac{1}{2}}}{p^{\frac{1}{2}}},\]
under $|S|\le p^3-p$.
Therefore, Theorem \ref{incidence} is better than Theorem \ref{thm: 3.8}.
\begin{proof}[Proof of Theorem \ref{thm: 3.8}]
     Since the identity matrix $I$ contributes at most $\min\{|A|, |B|\}$ incidences to $I(P,S)$, we may assume without loss generality that $I \not\in S$.
    For $\theta\in S$ and set $D \subset \Fbb_p^2 \times \Fbb_p^2$, by $i_D(\theta)$ we mean the number of pairs $(x, y)\in D$ such that $\theta y=x$. For $(x,y)\in P$, let $n(x,y)$ be the number of elements $(u,v) \in P$ such that $(u,v)=\lambda (x,y)$ for some $\lambda \in \Fbb_p^\ast$. Then  
\begin{align*}
    I(P, S)\le \sum_{\theta\in S}i_P(\theta) & = \sum_{\theta \in S}\sum_{(x,y) \in P , n(x,y)=1} \theta (x,y)+ \sum_{\theta \in S}\sum_{(x,y) \in P , n(x,y)>1} \theta (x,y)\\ 
    & =: I_1 +I_2
    % &\le |S|^{\frac{1}{2}}\left(\sum_{\theta \in S}i(\theta)^2\right)^{\frac{1}{2}} \\
    % & \ll \lv S\rv^{\frac{1}{2}} \l( \sum_{\theta \in S} \binom{i(\theta)}{2} +\lv S\rv \r)^{\frac{1}{2}} \\
    % & \le \lv S\rv^{\frac{1}{2}} \l( \lv \lo \l( x_1,y_1,x_2,y_2 \r) \in A\times B\times A \times B \colon \exists \theta \in S , \theta y_1 =x_1 ,\theta y_2 =x_2 \ro \rv +\lv S\rv \r)^{\frac{1}{2}}\\
    % & \le \lv S\rv^{\frac{1}{2}} \l( \lv \lo \l( x_1,y_1,x_2,y_2 \r) \in A\times B\times A \times B \colon x_1\cdot x_2^\perp =y_1\cdot y_2^\perp \ro \rv +\lv S\rv \r)^{\frac{1}{2}},
\end{align*} 
where $\theta (x,y)=1$ if $\theta y=x$, and $0$ ortherwise. Let $P'$ be the set of $(x, y)\in P$ such that $n(x, y)=1$. By the Cauchy--Schwarz inequality and Lemma \ref{lmmay3}
\begin{align*}
    I_1&=\sum_{\theta \in S } i_{P'}
    (\theta) \le |S|^{\frac{1}{2}}\left(\sum_{\theta \in S}i_{P'}(\theta)^2\right)^{\frac{1}{2}} \ll \lv S\rv^{\frac{1}{2}} \l( \sum_{\theta \in S} \binom{i_{P'}(\theta)}{2} +\lv S\rv \r)^{\frac{1}{2}} \\
    & \le \lv S\rv^{\frac{1}{2}} \l( \lv \lo \l( x_1,y_1,x_2,y_2 \r) \in A\times B\times A \times B \colon \exists \theta \in S , \theta y_1 =x_1 ,\theta y_2 =x_2 \ro \rv +\lv S\rv \r)^{\frac{1}{2}}\\
    & \le \lv S\rv^{\frac{1}{2}} \l( \lv \lo \l( x_1,y_1,x_2,y_2 \r) \in A\times B\times A \times B \colon x_1\cdot x_2^\perp =y_1\cdot y_2^\perp \ro \rv +\lv S\rv \r)^{\frac{1}{2}}\\
    & \ll \lv S\rv^{\frac{1}{2}} \l(  \frac{|A||B|}{p^{\frac{1}{2}}} + (pk|A||B|)^{\frac{1}{2}} + \lv S\rv^{\frac{1}{2}} \r) .
\end{align*}
We now consider $I_2$,
\begin{align*}
    I_2 &= \sum_{\theta \in S}\sum_{(x,y) \in P , n(x,y)>1} \theta (x,y)=\sum_{i=1}^{\lf \log_2 (k) \rf} \sum_{\theta\in S}\sum_{(x, y)\in P, 2^i \le n(x, y) < 2^{i+1}}\theta(x, y) .
\end{align*}
By using pigeonhole principle, there exists $\ell =2^{i_0} $ such that
\[ I_2 \lesssim \sum_{\theta\in S}\sum_{(x, y)\in P, \ell \le n(x, y)< 2\ell}\theta(x, y) . \]
Define $P''$ to be the set of representatives $(x, y)$ in $P$ such that $\ell \le n(x, y)< 2\ell$. Then, by the Cauchy--Schwarz inequality, Lemma \ref{lmmay3} and note that $k \ge \ell$ we have 
\begin{align*}
    I_2 & \lesssim \ell\sum_{\theta\in S}i_{P''}(\theta) \le \ell |S|^{\frac{1}{2}}\left(\sum_{\theta \in S}i_{P''}(\theta)^2\right)^{\frac{1}{2}}  \ll \lv S\rv^{\frac{1}{2}} \l( \sum_{\theta \in S}\ell^2 \binom{i_{P''}(\theta)}{2} + \ell^2\lv S\rv \r)^{\frac{1}{2}} \\
    & \le \lv S\rv^{\frac{1}{2}}  \l( \ell^2 \lv \lo  (x,y) \in (P'')^2\colon \exists \theta \in S , \theta \text{ is incident to $x$ and $y$} \ro \rv +\ell^2 \lv S\rv \r)^{\frac{1}{2}} \\
    & \le \lv S\rv^{\frac{1}{2}} \l( \lv \lo \l( x_1,y_1,x_2,y_2 \r) \in A\times B\times A \times B \colon \exists \theta \in S , \theta y_1 =x_1 ,\theta y_2 =x_2 \ro \rv +\ell^2 \lv S\rv \r)^{\frac{1}{2}}\\
    & \le \lv S\rv^{\frac{1}{2}} \l( \lv \lo \l( x_1,y_1,x_2,y_2 \r) \in A\times B\times A \times B \colon x_1\cdot x_2^\perp =y_1\cdot y_2^\perp \ro \rv +\ell^2 \lv S\rv \r)^{\frac{1}{2}}\\
    & \ll \lv S\rv^{\frac{1}{2}} \l(  \frac{|A||B|}{p^{\frac{1}{2}}} + (pk|A||B|)^{\frac{1}{2}} + k\lv S\rv^{\frac{1}{2}} \r) ,
\end{align*} 
% the fifth inequality follows from that for each $(x,y)\in (P'')^2$, there are $\ge \ell^2$ quadruples $(x_1,y_1,x_2,y_2) \in A\times B\times A \times B$ such that $(x_1,y_1) = \lambda_1 x, (x_2,y_2) =\lambda_2 y $, for some $ \lambda_1 ,\lambda_2 \in \Fbb_p^\ast$.
Putting these bounds together implies
\begin{align*}
    I(P,S) \lesssim \lv S\rv^{\frac{1}{2}} \l(  \frac{|A||B|}{p^{\frac{1}{2}}} + (pk|A||B|)^{\frac{1}{2}} + k\lv S\rv^{\frac{1}{2}} \r).
\end{align*}
as desired.
\end{proof}


%\newpage

% Compare the bound in Theorem \ref{incidence} and the bound in Theorem \ref{thm:3.7}: We have
% \begin{align*}
%     k^{\frac{1}{2}}|A|^{\frac{1}{2}}|S| < k^{\frac{1}{2}}p^{\frac{1}{2}}|S|^{\frac{1}{2}}|A|^\frac{1}{2}|B|^\frac{1}{2} & \Leftrightarrow |S| < p|B|, \\
%     |A|^{\frac{1}{2}}|B||S|^{3/4}p^{\frac{-1-\epsilon}{4}} < k^{\frac{1}{2}}p^{\frac{1}{2}}|S|^{\frac{1}{2}}|A|^\frac{1}{2}|B|^\frac{1}{2} & \Leftrightarrow |B|^2|S| < k^2p^{3+\epsilon}, \\
%     k^{\frac{1}{4}}p^{\frac{1-\epsilon}{4}}|A|^{\frac{1}{2}}|B|^\frac{1}{2}|S|^{\frac{3}{4}} < k^{\frac{1}{2}}p^{\frac{1}{2}}|S|^{\frac{1}{2}}|A|^\frac{1}{2}|B|^\frac{1}{2} & \Leftrightarrow |S| < kp^{1+\epsilon}.
% \end{align*}
% Therefore, if we have $|B|< k^{\frac{1}{2}}p$ and $|S|< \min \lo p|B|, kp^{1+\epsilon}\ro$ then the bound in Theorem \ref{thm:3.7} is better than the bound in Theorem \ref{incidence}.

% Compare with
% % the bound in Theorem \ref{thm:3.7} and the 
% bound in Theorem \ref{thm: 4.12}: 

% \begin{itemize}
%     \item When $|B|< k^{\frac{1}{2}}p $, we have
%     \begin{align*}
%     \frac{|A|^\frac{1}{2}|B||S|}{p}< k^{\frac{1}{2}}|A|^{\frac{1}{2}}|S|.
% \end{align*}
% So, the bound in Theorem \ref{thm: 4.12} is better than the bound in Theorem \ref{thm:3.7}.
%     \item When $k^{\frac{1}{2}}p \le |B|$, we have 
%     \begin{align*}
%     \frac{|A|^\frac{1}{2}|B||S|}{p}, k^{\frac{1}{4}}p^{\frac{1-\epsilon}{4}}|A|^{\frac{1}{2}}|B|^\frac{1}{2}|S|^{\frac{3}{4}} \le p^{\frac{-1-\epsilon}{4}}|A|^{\frac{1}{2}}|B||S|^{\frac{3}{4}} .
% \end{align*}
% \end{itemize}
% So, the bound in Theorem \ref{thm: 4.12} is better than the bound in Theorem \ref{thm:3.7}.


% \begin{align*}
%     |A|^{\frac{1}{2}}|B||S|^{3/4}p^{\frac{-1-\epsilon}{4}} < p^{\frac{2-\epsilon}{4}}|A|^{\frac{1}{2}}|B|^\frac{1}{2}|S|^{\frac{3}{4}} & \Leftrightarrow |B| < p^{\frac{3}{2}}, \\
%     k^{\frac{1}{2}}|A|^{\frac{1}{2}}|S| < p^{\frac{2-\epsilon}{4}}|A|^{\frac{1}{2}}|B|^\frac{1}{2}|S|^{\frac{3}{4}} & \Leftrightarrow |S| < p^{2-\epsilon}\frac{|B|^2}{k^2}.
% \end{align*}
% Therefore, when $k=o(p)$, if we have $|B|<p^\frac{3}{2}$ and $|S|< p^{2-\epsilon} \frac{|B|^2}{k^2}$ then the bound in Theorem \ref{thm:3.7} is better than the bound in Theorem \ref{thm: 4.12}.

% When $k \gg p$, we have $p^{\frac{2-\epsilon}{4}}|A|^{\frac{1}{2}}|B|^\frac{1}{2}|S|^{\frac{3}{4}} \ll k^{\frac{1}{4}}p^{\frac{1-\epsilon}{4}}|A|^{\frac{1}{2}}|B|^\frac{1}{2}|S|^{\frac{3}{4}}$ for all $A,B,S$ satisfied two theorems and
% \begin{align*}
%     \frac{|A|^\frac{1}{2}|B||S|}{p} < k^\frac{1}{2}|A|^{\frac{1}{2}}|S| & \Leftrightarrow |B|< k^{\frac{1}{2}}p.
%     % \frac{|A|^\frac{1}{2}|B||S|}{p} < |A|^{\frac{1}{2}}|B||S|^{3/4}p^{\frac{-1-\epsilon}{4}} &\Leftrightarrow |S| < p^{3-\epsilon}
% \end{align*}
% Therefore, when $k\gg p$, if we have $|B|<k^{\frac{1}{2}}p$ then the bound in Theorem \ref{thm: 4.12} is better than the bound in Theorem \ref{thm:3.7}.

% \begin{center}
% \begin{tabular}{ | m{3.2cm} | m{1cm}| m{2.5cm} | m{2.7cm} | m{2.5cm} | } 
%   \hline
%    & $|B|$ & $|B|< k^{\frac{1}{2}}p$ & $k^{\frac{1}{2}} p\le |B| < p^{\frac{3}{2}}$ & $p^{\frac{3}{2}} \le |B|$ \\ 
%   \hline
%   $|S|< p^{2-\epsilon}$ & cell5 & Theorem \ref{thm:3.7}  & Theorem \ref{thm:3.7} & Theorem \ref{thm: 4.12}\\ 
%   \hline
%   $p^{2-\epsilon}\le |S|<p^{2+\epsilon} $ & cell8 &   ???  & Theorem \ref{thm:3.7} & \\ 
%   \hline
%   $ p^{2+\epsilon} \le |S| $ & cell8 &    &  & Theorem \ref{incidence}???  \\
%   \hline
% \end{tabular}
% \end{center}

% \newpage

\subsection{Incidence bounds for small sets (Theorem \ref{incidence-energy'})}

%This section is devoted to prove an analog of Theorem \ref{incidence} for small sets. As mentioned in Subsection \ref{section2}, the incidence bound is sharp, so if one wishes to obtain improvements, some additional conditions are needed.  






%\begin{theorem}
%    Let $p$ be a prime and let $P =A \times B \subseteq \l( \Fbb_p^2 \times \Fbb_p^2 \r) \setminus \lo 0\ro $, $\lv B\rv \le p$. For $0< \gamma <\frac{3}{4}$, let $S$ be a symmetric subset of $\SL_2 \l( \Fbb_p \r)$ such that $p^\gamma < \lv S\rv < p^{3-2\gamma} $ and $\lv S\cap gH \rv < p^{\frac{-\gamma}{2}}\lv S\rv$ for any subgroup $H \subsetneq \SL_2 \l( \Fbb_p \r)$ and $g \in \SL_2 \l( \Fbb_p \r)$. Assume any line passing through the origin contains at most $k_1$ points from $B$ and $B$ determines at most $k_2$ distinct directions through the origin. 
%    Then, if $|B| \ge k_2k_1^{\frac{1}{2}}$, we have 
%    \[ I(P,S) \lesssim k_1^{\frac{1}{2}}|A|^{\frac{1}{2}}|S|+  k_1^{\frac{1}{2}}p^{\frac{-\epsilon}{4}}\lv B\rv^{\frac{1}{2}}|A|^{\frac{1}{2}}|S|^{\frac{3}{4}} + p^{\frac{-\epsilon}{4}} k_1^{\frac{1}{8}}\lv B\rv^{\frac{3}{4}}|A|^{\frac{1}{2}} |S|^{\frac{3}{4}},\]
%    if $|B| < k_2k_1^{\frac{1}{2}}$, we have 
%     \[ I(P,S) \lesssim k_1^{\frac{1}{2}}|A|^{\frac{1}{2}}|S|+  k_1^{\frac{1}{2}}p^{\frac{-\epsilon}{4}}\lv B\rv^{\frac{1}{2}}|A|^{\frac{1}{2}}|S|^{\frac{3}{4}} + k_1^{\frac{1}{4}}k_2^{\frac{1}{4}}p^{\frac{-\epsilon}{4}} \lv B\rv^{\frac{1}{2}} |A|^{\frac{1}{2}}|S|^{\frac{3}{4}},\]
%    where 
    
%    $\epsilon =\epsilon \l( \gamma \r)$ is a constant depending only on $\gamma$.
%\end{theorem}

% While the proof of Theorem \ref{incidence} uses the discrete Fourier analysis, for the case of small sets, our key ingredients are energy estimates. 



\subsubsection*{A skew dot-product energy estimate}
Given $B\subset \mathbb{F}_p^2$, this section is devoted to study the magnitude of the set 
\begin{equation}\label{ene-B}\{(x, y, u, v)\in B^4\colon x\cdot y^\perp=u\cdot v^\perp\}\end{equation}
when the size of $B$ is small. 

When the set $B$ is of large size, Lemma \ref{lmmay3} can be used to show that the number of such quadruples is almost the expected value $|B|^4/p$. However, when the size of $B$ is small, say, $|B|<p$, we have to deal with degenerate structures. 

More precisely, let $\ell$ be a line passing through the origin and $B$ be a subset of $\ell\setminus \{(0, 0)\}$, then the number of tuples $(x,y,u,v) \in B^4 $ such that $x\cdot y^\perp =u\cdot v^\perp $ is $\lv B\rv^4$. Note that if we remove the $'\perp'$ sign, then it will be at most $|B|^3$. 

This example shows that we need to have some conditions on the structures of $B$ so that a non-trivial estimate can be obtained. 


If $B$ has a few distinct directions through the origin, then, by following  Rudnev's argument in \cite[Section 3]{R.14} identically, one obtains
\begin{lemma}\label{lem: 6.12}
Let $B$ be a subset in $\mathbb{F}_p^2$ with $|B|\le p$. Assume any line passing through the origin contains at most $k_1$ points from $B$ and $B$ determines at most $k_2$ distinct directions through the origin. The number of quadruples $(x, y, u, v)\in B^4$ such that $x\cdot y^\perp=u\cdot v^\perp$ is at most $k_1^{\frac{1}{2}}|B|^3+k_1k_2|B|^2+k_1^2|B|^2$.
\end{lemma}

If we only assume an assumption on the number of points on a line passing through the origin, then we have the following lemma.

\begin{lemma}\label{lem: 6.2}
Let $B$ be a subset in $\mathbb{F}_p^2$ with $|B|\le p^{\frac{8}{15}}$. Assume any line passing through the origin contains at most $k$ points from $B$. The number of quadruples $(x, y, u, v)\in B^4$ such that $x\cdot y^\perp=u\cdot v^\perp$ is at most $k^{\frac{4}{15}}|B|^{\frac{748}{225}}+ k^2 \lv B\rv^2$.
\end{lemma}

\subsubsection*{Proof of Lemma \ref{lem: 6.2}}
The key ingredient in the proof of Lemma \ref{lem: 6.2} is the multi-set version of a point-line incidence bound due to Stevens and De Zeeuw in \cite{SdZ}.
\begin{theorem}[Stevens--de~Zeeuw, \cite{SdZ}]\label{frank}
Let $P$ be a point set and $L$ a set of lines in $\mathbb{F}_p^2$. If $|P|\ll p^{\frac{8}{5}}$, then the number of incidences between $P$ and $L$, denoted by $I(P, L)$, satisfies 
\[I(P, L)\ll |P|^{\frac{11}{15}}|L|^{\frac{11}{15}}+|P|+|L|.\]
\end{theorem}
\begin{theorem}\label{eqn:IPLMS}
    Let $P$ be a multi-set of points and $L$ be a multi-set of lines in $\mathbb{F}_p^2$. We denote the set of distinct points in $P$ by $\overline{P}$ and the set of distinct lines in $L$ by $\overline{L}$. For $p\in \overline{P}$ and $\ell\in \overline{L}$, let $m(p)$ and $m(\ell)$ be the multiplicity of $p$ and $\ell$, respectively. If 
    $|P|=\sum_{p\in \overline{P}}m(p)\ll p^{\frac{8}{5}}$, then
\begin{equation}\label{eqan:IPLMS}
    I(P, L)\lesssim |P|^{\frac{7}{15}}|L|^{\frac{7}{15}}\left(\sum_{p\in \overline{P}}m(p)^2\right)^{\frac{4}{15}}\left(\sum_{\ell\in \overline{L}}m(\ell)^2\right)^{\frac{4}{15}}+|P|+|L|.
\end{equation}
\end{theorem}
\begin{proof}
    Our argument to prove this theorem is similar to proof of \cite[Lemma 2.12]{LuPe}. 

Let $L_k$ be the set of lines in $\overline{L}$ of multiplicity $\sim 2^k$, and $P_k$ be the set of points in $\overline{P}$ of multiplicity $\sim 2^k$. Set 
\[Q_1:=\sum_{p\in \overline{P}}m(p)^2 \quad \text{and}\quad ~Q_2:=\sum_{\ell\in \overline{L}}m(\ell)^2.\]

Then, it is clear that 
\[\sum_{k}2^k|P_k|=|P|,\quad ~\sum_{k}2^{2k}|P_k|=Q_1,\]
and
\[\sum_{k}2^k|L_k|=|L|,\quad ~\sum_{k}2^{2k}|L_k|=Q_2.\]
Thus, we have 
\[|P_k|\le \min \left\lbrace\frac{|P|}{2^k}, \frac{Q_1}{2^{2k}}\right\rbrace \quad \text{and} \quad |L_k|\le \min \left\lbrace\frac{|L|}{2^k}, \frac{Q_2}{2^{2k}}\right\rbrace.\]

We now observe 
\begin{align*}
    I(P, L)&=\sum_{i, j}I(P_i, L_j)=\sum_{2^j<Q_2/|L|}2^jI(P, L_j)+\sum_{2^j\ge Q_2/|L|}I(P, L_j)\\
    &=\sum_{\substack{2^i<Q_1/|P|\\2^j<Q_2/|L|}}2^i2^jI(P_i, L_j)+\sum_{\substack{2^i\ge Q_1/|P|\\2^j<Q_2/|L|}}2^i2^jI(P_i, L_j)\\&+\sum_{\substack{2^i<Q_1/|P|\\2^j\ge Q_2/|L|}}2^i2^jI(P_i, L_j)+\sum_{\substack{2^i\ge Q_1/|P|\\2^j\ge Q_2/|L|}}2^i2^jI(P_i, L_j)\\
    &=: I+II+III+IV.\\
\end{align*}
{\bf Bounding $I$:} Since $|P_i|\le |P|/2^i$ and $|L_j|\le |L|/2^j$, by Theorem \ref{frank}, we obtain 
\begin{align*}
I&\ll \sum_{\substack{2^i<Q_1/|P|\\
2^j<Q_2/|L|}}2^{i+j}\left(\frac{|P||L|}{2^{i+j}}\right)^{\frac{11}{15}}+|P|+|L|=\sum_{\substack{2^i<Q_1/|P|\\
2^j<Q_2/|L|}}2^{\frac{4(i+j)}{15}}(|P||L|)^{\frac{11}{15}}+|P|+|L|\\
&\lesssim |P|^{\frac{7}{15}}|L|^{\frac{7}{15}}(Q_1Q_2)^{\frac{4}{15}}+|P|+|L|.
\end{align*}

{\bf Bounding $II$:} Using $|P_i|\le Q_\frac{1}{2}^{2i}$ and $|L_j|\le |L|/2^j$, Theorem \ref{frank} implies
\begin{align*}
    II&:=\sum_{\substack{2^i\ge Q_1/|P|\\2^j<Q_2/|L|}}2^i2^jI(P_i, L_j)\le \sum_{\substack{2^i\ge Q_1/|P|\\2^j<Q_2/|L|}}2^i2^j \left(\frac{Q_1}{2^{2i}}\right)^{\frac{11}{15}}\left(\frac{|L|}{2^j}\right)^{\frac{11}{15}}+|P|+|L|\\
    &\lesssim |P|^{\frac{7}{15}}|L|^{\frac{7}{15}}(Q_1Q_2)^{\frac{4}{15}}+|P|+|L|.
\end{align*}

{\bf Bounding $III, IV$:} Similarly, we also have 
\[III, IV\lesssim |P|^{\frac{7}{15}}|L|^{\frac{7}{15}}(Q_1Q_2)^{\frac{4}{15}}+|P|+|L|.\]

This concludes the proof of the theorem.
\end{proof}
With this incidence bound, we are ready to prove Theorem \ref{lem: 6.2}.
\begin{proof}[Proof of Theorem \ref{lem: 6.2}]
We have
\begin{align*}
    \lv \{(x,y,u,v) \in B^4 \colon x\cdot y^\perp =u\cdot v^\perp \}\rv & = \lv \{(x,y,u,v) \in B^4 \colon x\cdot y^\perp =u\cdot v^\perp \ne 0\}\rv \\&+\lv \{(x,y,u,v) \in B^4 \colon x\cdot y^\perp =u\cdot v^\perp =0\}\rv \\
    & =: I +II
\end{align*}
Regarding $I$. For $a, u, v\in B$, $u\cdot v^\perp \ne 0$ we define $L_{a, u, v}$ to be the multi-set of lines of the form $x\cdot a^\perp=u\cdot v^\perp$. Let $L=\bigcup_{a, u, v\in B}L_{a, u, v}$. It is clear that the number of such quadruples is at most $I(B, L)$. We note that $|L|\le |B|^3$. By applying Theorem \ref{frank}, we observe that for each line in $L$, its multiplicity is at most $\ll k|B|^{\frac{22}{15}}$. For each $l \in L$, let $m(l)$ be the multiplicity $l$, we have
\[\sum_{\ell}m(\ell)=|L|\le |B|^3,\]
and 
\[\sum_{\ell}m(\ell)^2\le \max_{\ell}m(\ell)\cdot |B|^3\ll k|B|^{\frac{22}{15}}\cdot |B|^3.\]
It follows from Theorem \ref{eqn:IPLMS} for $B$ and $L$ that
\[ I(B,L) \lesssim \lv B\rv^{\frac{7}{15}} \lv B\rv^{\frac{21}{15}}\lv B\rv^{\frac{4}{15}}\l(  k\lv B\rv^{\frac{67}{15}} \r)^{\frac{4}{15}} +\lv B\rv +\lv L\rv \ll k^{\frac{4}{15}} \lv B\rv^{\frac{748}{225}}. \]

Regarding $II$. For each $x,u\in B$, there are at most $k$ points $y$ and $k$ points $v$ such that $x\cdot y^\perp =u\cdot v^\perp =0$. So 
\[ II \le k^2\lv B\rv^{2}. \]
This completes the proof.
\end{proof}

\subsubsection*{Proof of Theorem \ref{incidence-energy'}}
We follow the proof of Theorem \ref{thm:3.7} with Lemma \ref{lem: 6.12} and Lemma \ref{lem: 6.2} in place of Lemma \ref{lmmay3}, one has two following bounds
\begin{align*} I(P,S) &\lesssim  k_1^{\frac{1}{2}}|A|^{\frac{1}{2}}|S|+  \frac{k_1^{\frac{1}{2}}\lv B\rv^{\frac{1}{2}}|A|^{\frac{1}{2}}|S|^{\frac{3}{4}}}{p^{\frac{\epsilon}{4}}} + \frac{k_1^{\frac{1}{8}}\lv B\rv^{\frac{3}{4}}|A|^{\frac{1}{2}} |S|^{\frac{3}{4}}}{p^{\frac{\epsilon}{4}}}+\frac{k_1^{\frac{1}{4}}k_2^{\frac{1}{4}}\lv B\rv^{\frac{1}{2}} |A|^{\frac{1}{2}}|S|^{\frac{3}{4}}}{p^{\frac{\epsilon}{4}}},
\end{align*}
and 
\[ I(P,S) \lesssim  k_1^{\frac{1}{2}}|A|^{\frac{1}{2}}|S|+  \frac{k_1^{\frac{1}{2}}\lv B\rv^{\frac{1}{2}}|A|^{\frac{1}{2}}|S|^{\frac{3}{4}}}{p^{\frac{\epsilon}{4}}}+\frac{k_1^{\frac{1}{15}}|B|^{\frac{187}{225}} |A|^{\frac{1}{2}} |S|^{\frac{3}{4}}}{p^{\frac{\epsilon}{4}}}.\]
This completes the proof. $\square$

For applications, a comparison of these two incidence bounds, with $|B|\le p$, is provided as follows. 
\begin{enumerate}
    \item If $|B|\ge k_2k_1^{\frac{1}{2}}$, then
%\[ \frac{k_1^{\frac{1}{4}}k_2^{\frac{1}{4}}\lv B\rv^{\frac{1}{2}} |A|^{\frac{1}{2}}|S|^{\frac{3}{4}}}{p^{\frac{\epsilon}{4}}} \le \frac{k_1^{\frac{1}{8}}\lv B\rv^{\frac{3}{4}}|A|^{\frac{1}{2}} |S|^{\frac{3}{4}}}{p^{\frac{\epsilon}{4}}} \le  \frac{k_1^{\frac{1}{15}}|B|^{\frac{187}{225}} |A|^{\frac{1}{2}} |S|^{\frac{3}{4}}}{p^{\frac{\epsilon}{4}}} .\]
%So, we have
\[ I(P,S) \lesssim  k_1^{\frac{1}{2}}|A|^{\frac{1}{2}}|S|+  \frac{k_1^{\frac{1}{2}}\lv B\rv^{\frac{1}{2}}|A|^{\frac{1}{2}}|S|^{\frac{3}{4}}}{p^{\frac{\epsilon}{4}}}+ \frac{k_1^{\frac{1}{8}}\lv B\rv^{\frac{3}{4}}|A|^{\frac{1}{2}} |S|^{\frac{3}{4}}}{p^{\frac{\epsilon}{4}}}. \]
\item If $k_1^{\frac{165}{298}}k_2^{\frac{225}{298}} \le |B|< k_2k_1^{\frac{1}{2}}$, then
%\[ \frac{k_1^{\frac{1}{8}}\lv B\rv^{\frac{3}{4}}|A|^{\frac{1}{2}} |S|^{\frac{3}{4}}}{p^{\frac{\epsilon}{4}}} \le \frac{k_1^{\frac{1}{4}}k_2^{\frac{1}{4}}\lv B\rv^{\frac{1}{2}} |A|^{\frac{1}{2}}|S|^{\frac{3}{4}}}{p^{\frac{\epsilon}{4}}} \le \frac{k_1^{\frac{1}{15}}|B|^{\frac{187}{225}} |A|^{\frac{1}{2}} |S|^{\frac{3}{4}}}{p^{\frac{\epsilon}{4}}}. \]
%So, we have
\[ I(P,S) \lesssim  k_1^{\frac{1}{2}}|A|^{\frac{1}{2}}|S|+  \frac{k_1^{\frac{1}{2}}\lv B\rv^{\frac{1}{2}}|A|^{\frac{1}{2}}|S|^{\frac{3}{4}}}{p^{\frac{\epsilon}{4}}}+ \frac{k_1^{\frac{1}{4}}k_2^{\frac{1}{4}}\lv B\rv^{\frac{1}{2}} |A|^{\frac{1}{2}}|S|^{\frac{3}{4}}}{p^{\frac{\epsilon}{4}}}. \]
\item If $|B| <k_1^{\frac{165}{298}}k_2^{\frac{225}{298}} $, then
%\[ \frac{k_1^{\frac{1}{8}}\lv B\rv^{\frac{3}{4}}|A|^{\frac{1}{2}} |S|^{\frac{3}{4}}}{p^{\frac{\epsilon}{4}}},\frac{k_1^{\frac{1}{15}}|B|^{\frac{187}{225}} |A|^{\frac{1}{2}} |S|^{\frac{3}{4}}}{p^{\frac{\epsilon}{4}}} \le \frac{k_1^{\frac{1}{4}}k_2^{\frac{1}{4}}\lv B\rv^{\frac{1}{2}} |A|^{\frac{1}{2}}|S|^{\frac{3}{4}}}{p^{\frac{\epsilon}{4}}} . \]
%So, we have 
\[ I(P,S) \lesssim  k_1^{\frac{1}{2}}|A|^{\frac{1}{2}}|S|+  \frac{k_1^{\frac{1}{2}}\lv B\rv^{\frac{1}{2}}|A|^{\frac{1}{2}}|S|^{\frac{3}{4}}}{p^{\frac{\epsilon}{4}}}+\frac{k_1^{\frac{1}{15}}|B|^{\frac{187}{225}} |A|^{\frac{1}{2}} |S|^{\frac{3}{4}}}{p^{\frac{\epsilon}{4}}}.\]
\end{enumerate}

\subsection{Alternative approach with relaxed conditions}
This section is devoted to prove the following analog, which offers a bound which is meaningful when the size of $S$ is large compared to the sizes of $A$ and $B$.


\begin{theorem}\label{thm: 6.2}
    Let $p$ be a prime and $P=A\times B \subseteq \l( \Fbb_p^2\times \Fbb_p^2 \r) \setminus \{ 0 \}$ with $\lv B\rv \le \lv A\rv \le p^{\frac{8}{15}}$ and any lines passing through the origin contains at most $k$ points from $B$. Then, we have 
    % \[ I(P,S) \ll \lv S\rv^{\frac{1}{2}} \l(  \min \lo k^{\frac{2}{15}}\lv B\rv^{\frac{11}{15}} \lv A\rv^{\frac{209}{225}} ,\lv A\rv^{\frac{13}{15}}\lv B\rv^{\frac{209}{225}} \ro +  k^{\frac{1}{2}}\lv A\rv \lv B\rv^{\frac{1}{2}}+ k\lv S\rv^{\frac{1}{2}} \r) .\]
    \[ I(P,S) \lesssim k^{\frac{2}{15}}\lv A\rv^{\frac{11}{15}}\lv B\rv^{\frac{209}{225}}  |S|^{\frac{1}{2}}  +k|A|^{\frac{1}{2}}|B|^{\frac{1}{2}}|S|^{\frac{1}{2}} +  k\lv S\rv .\]
\end{theorem}

% Theorem \ref{incidence-energy} is proved by following the proof of Theorem \ref{thm:3.7} and the next lemma. 





Theorem \ref{thm: 6.2} is proved by following the proof of Theorem \ref{thm: 3.8} and the next lemma, whose proof is identical to that of Lemma \ref{lem: 6.2}.


\begin{lemma}\label{lem: quad}
    Let $A,B \subset \Fbb_p^2$ with $\lv B\rv \le \lv A \rv \le p^{\frac{8}{15}}$. Assume any line passing through the origin contains at most $k$ points from $B$ and at most $k$ points from $A$. The number of quadruples $(x,y,u,v) \in A\times A \times B \times B $ such that $x\cdot y^\perp =u\cdot v^\perp$ is at most $  k^{\frac{4}{15}}\lv A\rv^{\frac{22}{15}}\lv B\rv^{\frac{418}{225}}  + k^2\lv A\rv\lv B\rv$. 
\end{lemma}

%     We have
% \begin{align*}
%     &  \lv (x,y,u,v) \in A\times A \times B \times B \colon x\cdot y^\perp =u\cdot v^\perp \rv \\
%     & = \lv (x,y,u,v) \in A\times A \times B \times B \colon x\cdot y^\perp =u\cdot v^\perp \ne 0\rv \\&+\lv (x,y,u,v) \in A\times A \times B \times B \colon x\cdot y^\perp =u\cdot v^\perp =0\rv \\
%     & =: I +II
% \end{align*}
% Regarding $I$. There exist two way to estimate $I$. On the one hand, for $(a, u, v)\in A\times B^2$, $u\cdot v^\perp \ne 0$ we define $L_{a, u, v}$ to be the multi-set of lines of the form $x\cdot a^\perp=u\cdot v^\perp$. Let $L=\bigcup_{(a, u, v)\in A\times B^2}L_{a, u, v}$. It is clear that the number of such quadruples is at most $I(A, L)$. We note that $|L|\le |A||B|^2$. By applying Theorem \ref{frank}, we observe that for each line in $L$, its multiplicity is at most $\ll k|B|^{\frac{22}{15}}$. For each $l \in L$, let $m(l)$ be the multiplicity $l$, we have
% \[\sum_{\ell}m(\ell)=|L|\le |A||B|^2,\]
% and 
% \[\sum_{\ell}m(\ell)^2\le \max_{\ell}m(\ell)\cdot |A||B|^2\ll k|B|^{\frac{22}{15}}\cdot |A||B|^2.\]
% It follows from Theorem \ref{eqn:IPLMS} for $A$ and $L$ that
% \[ I(A,L) \ll \lv A\rv^{\frac{7}{15}} \lv A\rv^{\frac{7}{15}}|B|^{\frac{14}{15}}\lv A\rv^{\frac{4}{15}}\l(  k|A|\lv B\rv^{\frac{52}{15}} \r)^{\frac{4}{15}} +\lv A\rv +\lv L\rv \ll k^{\frac{4}{15}}|A|^{\frac{22}{15}} \lv B\rv^{\frac{418}{225}}. \]

% On the other hand, for $(b, x, y)\in B\times A^2$, $x\cdot y^\perp \ne 0$ we define $L_{b, x, y}$ to be the multi-set of lines of the form $z\cdot b^\perp=x\cdot y^\perp$. Let $L_1=\bigcup_{(b, x, y)\in B\times A^2}L_{b, x, y}$. It is clear that the number of such quadruples is at most $I(B, L_1)$. We note that $|L_1|\le |A|^2|B|$. By applying Theorem \ref{frank}, we observe that for each line in $L$, its multiplicity is at most $\ll k|A|^{\frac{22}{15}}$. For each $l \in L_1$, let $m(l)$ be the multiplicity $l$, we have
% \[\sum_{\ell}m(\ell)=|L_1|\le |A|^2|B|,\]
% and 
% \[\sum_{\ell}m(\ell)^2\le \max_{\ell}m(\ell)\cdot |A|^2|B|\ll k|A|^{\frac{22}{15}}\cdot |A|^2|B|.\]
% It follows from Theorem \ref{eqn:IPLMS} for $B$ and $L_1$ that
% \[ I(B,L_1) \ll \lv B\rv^{\frac{7}{15}} \lv A\rv^{\frac{14}{15}}|B|^{\frac{7}{15}}\lv B\rv^{\frac{4}{15}}\l( k|B|  |A|^{\frac{52}{15}} \r)^{\frac{4}{15}} +\lv B\rv +\lv L_1\rv \ll k^{\frac{4}{15}}|B|^{\frac{22}{15}}|A|^{\frac{418}{225}}. \]

% Regarding $II$. For each $(x,u)\in A\times B$ there are at most $k$ points $y$ and $k$ points $v$ respectively such that $x,y$ and $u,v$ contain in same lines through the origin. Therefore, there are at most $k$ points $y$ and $k$ points $v$ such that $x\cdot y^\perp =u\cdot v^\perp =0$. Then. we have 
% \[ II \le k^2\lv A\rv \lv B\rv. \]
% This completes the proof.

\section{Proof of Theorems \ref{thm2} -- \ref{thm3'}}
%We now use Theorem \ref{incidence} to prove Theorems \ref{thm2} and \ref{thm1}.
\begin{proof}[Proof of Theorem \ref{thm2}]
On the one hand, we apply Theorem \ref{incidence} for $S$ and $P= S(E) \times E$ to obtain
\begin{align}\label{eq: 1}
    I(S(E)\times E,S) \ll \frac{\lv S(E)\rv \lv E\rv \lv S\rv}{p^2} + p\sqrt{\lv S(E)\rv \lv E\rv \lv S\rv} + \lv S\rv.
\end{align}
On the other hand, for each $x\in E$ and $\theta \in S$, there exists unique $y\in S(E)$ such that $(y,x)$ is incident to $\theta$. So, we obtain
 \[|E||S|=I(S(E)\times E,S) \ll \frac{\lv S(E)\rv \lv E\rv \lv S\rv}{p^2} + p\sqrt{\lv S(E)\rv \lv E\rv \lv S\rv} + \lv S\rv.\]
 Solving this inequality, the theorem follows. 
\end{proof}
% Theorem \ref{thm11} and 
Theorem \ref{thm555} will follow from the two following results. 
\begin{theorem}\label{thm1}
    Let $E\subset \mathbb{F}_p^2$ and $S\subset SL_2(\mathbb{F}_p)$.
\begin{itemize}
    \item[a.] If  $|E|\ge 4p$ and $|S|\gg p^2$, then there exists $x\in E$ such that $|S(E-x)|\gg p^2$. 
    \item[b.] If any line through the origin contains at most $k$ points from $E$, then 
     \[|S(E)|\gg \min \left\lbrace p^2, ~\frac{|S||E|}{pk} \right\rbrace.\]
\end{itemize}
\end{theorem}
Theorem \ref{thm1} (a) is not a part of Theorem \ref{thm555}, but we find it to be of independent interest. 

\begin{theorem}\label{thm5}
For any $\gamma\in (0, 1)$, there exists $\epsilon=\epsilon(\gamma)>0$ such that the following holds.

     Let $p$ be a sufficiently large prime. Let $E \subseteq \Fbb_p^2  \setminus \lo (0, 0)\ro$, and let $S$ be a symmetric subset of $\SL_2 (\Fbb_p )$ such that $p^\gamma < |S| < p^{3-2\gamma } $ and $ |S \cap gH| < p^{\frac{-\gamma}{2}}|S|$ for any subgroup $H \subsetneq \SL_2 (\Fbb_p) $ and $g\in \SL_2(\Fbb_p)$. Then, we have 
    \[ |S(E)| \gg \min \lo p^2 , \frac{|S|^{\frac{1}{2}}|E|}{p^{1-\frac{\epsilon}{2}}} \ro .\] Moreover,  if any line passing through the origin contains at most $k$ points from $E$, we have
    \[ |S(E)| \gg \min \lo p^2 ,\frac{|S|^{\frac{1}{2}}|E|}{p^{\frac{1-\epsilon}{2}}k^{\frac{1}{2}}} \ro .\]
\end{theorem}

To prove Theorem \ref{thm1}(a), we recall the following result from \cite{L.P.V.18}. 

\begin{lemma}[Corollary 10, \cite{L.P.V.18}]\label{exist}
    Let $E\subset \mathbb{F}_p^2$ with $|E|\ge 4p$. Then there exists a point $x\in E$ such that there are at least $p/2$ lines passing through $x$ and each line contains at least one other point from $E$. 
\end{lemma}


\begin{proof}[Proof of Theorem \ref{thm1}]\,
\medskip

{\bf Part a.}
    By Lemma \ref{exist}, there exists a point $x\in E$ such that there are at least $p/2$ lines passing through $x$ and each line contains at least one other point from $E$. From each of these lines, we pick one point which is different from $x$ and let $E'$ be the set of those points. Then we have $|E'|\ge p/2$. We observe that 
    \[|S(E-x)|\ge |S(E'-x)|.\]
    Note that $E'-x$ is a translation of $E'$ by $x$. So, any line passing through the origin contains at most one point from $E'-x$. Set $E''=E'-x$. We now estimate the size of $S(E'')$ from below. 

Set $P=E''\times S(E'')$ and $P'=\{\lambda\cdot u\colon u\in P, \lambda\ne 0\}$. We have $|P'|=(p-1)|P|$.
    Note that $I(P, S)=|E''||S|=\frac{1}{p-1}I(P', S)$. 
By Theorem \ref{incidence}, we have     
\[I(P', S)\ll \frac{|P'||S|}{p^2}+p\sqrt{|P'||S|} + \lv S\rv .\]
Putting the upper and lower bounds together, the theorem follows. 

{\bf Part b.} The proof is almost the same, we just need to partition the set $E$ into at most $k$ subsets $E_i$ such that each has the same structure as the set $E''$ in the Part a. So we omit the details.
\end{proof}

Theorem \ref{thm5}, Theorem \ref{thm345}, and Theorem \ref{thm3'} are proved by following the proof of Theorem \ref{thm2}, but we use Theorems \ref{thm: 4.12} and \ref{incidence-energy'} (1) in place of Theorem \ref{incidence}.

\begin{remark}
Theorem \ref{incidence-energy'} (2) implies the following bound
     \[ \lv S(E)\rv \gtrsim \min \lo \frac{\lv E\rv^{0.337} \lv S\rv^{\frac{1}{2}}p^{\epsilon /2}}{k_1^{\frac{2}{15}}}, \frac{\lv E\rv \lv S\rv^{\frac{1}{2}}p^{\epsilon /2}}{k_1}, \frac{\lv E\rv^2}{k_1}  \ro,\]
     which improves Theorem \ref{thm3'} in the range $|E|\le \min \lo p^{\frac{8}{15}} , k_1^{0.553}k_2^{0.755} \ro$. 
\end{remark}

% \begin{proof}[Proof of Theorem \ref{thm3'}]
%     Let $P= S(E)\times E$. Note that $I(P,S)=\lv E\rv \lv S\rv$. By using Theorem \ref{incidence-energy'}, 
%     \begin{itemize}
%         \item[(1)] when $|B|< k_1^{\frac{165}{298}}k_2^{\frac{225}{298}}$, we have
%     \[ \lv E\rv \lv S\rv = I(S(E) \times E ,S) \lesssim \lv S(E)\rv^{\frac{1}{2}} \l( k_1^{\frac{1}{2}}|S|+\l( k_1^{\frac{1}{2}}\lv E\rv^{\frac{1}{2}}+ k_1^{1/15}|E|^{187/225}\r) |S|^{3/4}p^{-\epsilon/4}\r) ,\]
%         \item[(2)] when $|B|\ge k_1^{\frac{165}{298}}k_2^{\frac{225}{298}}$, we have
%         \[  I(P,S) \lesssim k_1^{\frac{1}{2}}|A|^{\frac{1}{2}}|S|+  \frac{k_1^{\frac{1}{2}}\lv B\rv^{\frac{1}{2}}|A|^{\frac{1}{2}}|S|^{\frac{3}{4}}}{p^{\frac{\epsilon}{4}}} + \frac{k_1^{\frac{1}{8}}\lv B\rv^{\frac{3}{4}}|A|^{\frac{1}{2}} |S|^{\frac{3}{4}}}{p^{\frac{\epsilon}{4}}} + \frac{k_1^{\frac{1}{4}}k_2^{\frac{1}{4}}\lv B\rv^{\frac{1}{2}} |A|^{\frac{1}{2}}|S|^{\frac{3}{4}}}{p^{\frac{\epsilon}{4}}},\]
%     \end{itemize}
%     where $\epsilon =\epsilon (\gamma)$ is a constant depending only on $\gamma$.
%     Therefore, 
%      \[ \lv S(E)\rv \gtrsim \min \lo  \max \lo \min \lo \frac{\lv E\rv^{\frac{1}{2}} \lv S\rv^{\frac{1}{2}}p^{\epsilon /2}}{k_1^{1/4}}, \frac{\lv E\rv \lv S\rv^{\frac{1}{2}}p^{\epsilon /2}}{k_1^{\frac{1}{2}}k_2^{\frac{1}{2}}} \ro ,\frac{\lv E\rv^{76/225} \lv S\rv^{\frac{1}{2}}p^{\epsilon /2}}{k_1^{\frac{2}{15}}}   \ro , \frac{\lv E\rv \lv S\rv^{\frac{1}{2}}p^{\epsilon /2}}{k_1}, \frac{\lv E\rv^2}{k_1} \ro .\] 
% \end{proof}

% In the above argument, in stead of using Theorem \ref{incidence-energy}, 
% one can use Lemma \ref{incidence-energy'} and obtain Theorem \ref{thm3'}.




\section{Incidence structures spanned by \texorpdfstring{$\mathbb{H}_1(\mathbb{F}_p)$}{H1(Fp)}}
As in the case of the special linear group, we define an incidence structure as follows. We say the point $(x, y)\in \mathbb{F}_p^3\times \mathbb{F}_p^3$ is incident to the matrix $\theta\in \mathbb{H}_1(\Fbb_p)$ if $\theta y=x$. 

\begin{theorem}\label{H-inci}Let $\epsilon\ge 0$.
    Let $X\subset \mathbb{H}_1(\Fbb_p)$ and $P=A\times B\subset \mathbb{F}_p^6$. Assume that all points in $A$ and $B$ have the third coordinate non-zero, and for each $(y, z)\in \mathbb{F}_p^2$, we have $|\pi_{23}^{-1}(y, z)\cap B|\le p^{1-\epsilon}$, then the number of incidences between $X$ and $P$ satisfies
    \[\left\vert I(P, X)- \frac{|P||X|}{p^3} \right\vert\le p^{\frac{3-\epsilon}{2}}|P|^{\frac{1}{2}}|X|^{\frac{1}{2}}.\]
\end{theorem}

%With this theorem in hand, one can prove Theorem \ref{main-H} easily. To see this, we set $A=E$ and $B=X(E)$. We now estimate the number of incidences between $A\times B$ and $X$. It is clear that $I(A\times B, X)=|X||A|$. Applying Theorem \ref{H-inci} on the number of incidences between $A\times B$ and $X$, we obtain 
%\[I(A\times B, X)\le \frac{|A||B||X|}{p^3}+p^{\frac{3-\epsilon}{2}}|X|^{\frac{1}{2}}|A|^{\frac{1}{2}}|B|^{\frac{1}{2}}.\]
%Putting the lower and upper bounds together and solving the inequality give us the desired bound. 

We now turn our attention to proving Theorem \ref{H-inci}.
\begin{proof}[Proof of Theorem \ref{H-inci}]
By repeating the argument as in the case of $SL_2(\mathbb{F}_p)$, we have 
    \begin{align*}
        I(P,X) & = \sum_{\substack{(x,y )\in P \\ \theta \in X}} 1_{\theta x=y } =\frac{1}{p^3} \sum_{m\in \Fbb_p^3} \sum_{\substack{(x,y)\in P, \\ \theta \in X}} \chi \l( m\cdot \l( \theta y-x \r) \r) \\
        & = \frac{\lv P\rv \lv X\rv}{p^3} + \frac{1}{p^3} \sum_{ m\in \Fbb_p^3 \setminus \lo (0,0, 0)\ro } \sum_{\substack{ (x,y)\in P ,\\ \theta \in X }} \chi \l( m \cdot \l( \theta y -x \r) \r) \\
        & = \frac{\lv P\rv \lv X\rv}{p^3} + p^3\sum_{m\ne (0, 0, 0)} \sum_{\theta \in X} \widehat{P} \l( -m,\theta^t m\r).
    \end{align*}
By using the Cauchy--Schwarz inequality, one has 
\begin{align*}
    &\sum_{\theta\in X}\sum_{m\ne (0, 0, 0)}\widehat{P}(-m, \theta^tm)\le |X|^{\frac{1}{2}}\left(\sum_{\theta\in \mathbb{H}_1(\Fbb_p)}\sum_{m_1, m_2\ne (0, 0, 0)}\widehat{P}(-m_1, \theta^tm_1)\overline{\widehat{P}(-m_2, \theta^tm_2)}\right)^{\frac{1}{2}}.
\end{align*}
We now observe 
\begin{align*}
    &\sum_{\theta\in \mathbb{H}_1(\Fbb_p)}\sum_{m_1, m_2\ne (0, 0, 0)}\widehat{P}(-m_1, \theta^tm_1)\overline{\widehat{P}(-m_2, \theta^tm_2)}\\&=\frac{1}{p^{12}}\sum_{\theta}\sum_{m_1, m_2\ne (0, 0, 0)}\sum_{(x_1, y_1), (x_2, y_2)}P(x_1, y_1)P(x_2, y_2)\chi(-m_1x_1+\theta^tm_1y_1)\chi(m_2x_2-\theta^tm_2y_2)\\
    &=\sum_{m_1, m_2}-\sum_{m_1=(0, 0, 0), m_2\ne (0, 0, 0)}-\sum_{m_1\ne (0, 0, 0), m_2=(0, 0, 0)}-\sum_{m_1=m_2=(0, 0, 0)}\\
    &=I-II-III-IV.
\end{align*}
We now estimate each term separately. 

Regarding $IV$, it is clear that it is equal to $|P|^2/p^9$. The terms $II$ and $III$ are the same, and we can see that 
\[II=\frac{|P|}{p^{12}}\sum_{\theta}\sum_{m_2\ne (0, 0, 0)}\sum_{(x_2, y_2)}P(x_2, y_2)\chi(m_2(x_2-\theta y_2)),\]
which can be written as 
\[\frac{|P|(p^2-1)}{p^{12}}\sum_{\theta}\sum_{(x_2, y_2)\in P, x_2=\theta y_2} 1-\frac{|P|}{p^{12}}\sum_{\theta}\sum_{(x_2, y_2)\in P, x_2\ne\theta y_2}1.\]
So, $-II\le \frac{|P|^2}{p^9}$.
Regarding $I$,
\begin{align*}
    &I=\frac{1}{p^{12}}\sum_{\theta}\sum_{m_1, m_2\in \mathbb{F}_p^3}\sum_{(x_1, y_1), (x_2, y_2)}P(x_1, y_1)P(x_2, y_2)\chi(m_1(\theta y_1-x_1))\chi(m_2(\theta y_2-x_2))\\
    &=\frac{1}{p^6}\sum_{\theta}\sum_{(x_1, y_1), (x_2, y_2)}P(x_1, y_1)P(x_2, y_2)1_{\theta y_1=x_1}1_{\theta y_2=x_2}.\\
\end{align*}
To proceed further, we need to count the number of quadruples $(x, y, z)\in B, (x', y', z')\in A, (u, v, w)\in B, (u', v', w')\in A$ such that there exists $\theta=[a, b, c]\in \mathbb{H}_1(\Fbb_p)$ and $\theta (x, y, z)=(x', y', z')$, ~$\theta(u, v, w)=(u', v', w')$.

For such quadruples, one has 
\[yw'+zv'=y'w+vz', ~z(u'-u)+w(x-x')=a(vz-wy), ~z=z', w=w'.\]
From here, as in the case of the special linear group, we need to make use of a number of preliminary lemmas, which will be proved later.
\begin{proposition}\label{propo2.1}
Given $A, B\subset \mathbb{F}_p^3$, let $N$ be the number of quadruples $(x, y, z)\in A, (x', y', z')\in B, (u, v, w)\in A, (u', v', w')\in B$ such that 
\[yw'+zv'=y'w+vz', ~z=z', w=w'.\]
Suppose in addition that for each $(y, z)\in \mathbb{F}_p^2$, we have $\pi_{23}^{-1}(y, z)\cap A$ or $\pi_{23}^{-1}(y, z)\cap B$ are of sizes at most $p^{1-\epsilon}$, then we have 
\[N\le \frac{|A|^2|B|^2}{p}+p^{3-2\epsilon}|A||B|.\]
\end{proposition}
\begin{proposition}\label{propo2.2}
    Given $A, B\subset \mathbb{F}_p^3$, let $N'$ be the number of quadruples $(x, y, z)\in B, (x', y', z')\in A, (u, v, w)\in B, (u', v', w')\in A$ such that 
\[yw'+zv'=y'w+vz', ~vz-wy=0, ~zu'+xw'-z'u-x'w=0, ~w=w',~z=z'.\]
Then $N'\le p|A||B|$.
\end{proposition}
\begin{lemma}\label{lm23}
    Given $(x, y, z), (x', y', z'), (u, v, w), (u', v', w')$ in $\mathbb{F}_p^3$ with $w\ne 0$, $z\ne 0$, and $yw'+zv'=y'w+z'v$. Let $N''$ be the number of matrices $\theta=[a, b, c]\in \mathbb{H}_1(\Fbb_p)$ such that 
\[\theta(x, y, z)=(x', y', z')\]
and 
\[\theta(u, v, w)=(u', v', w').\]
\begin{itemize}
    \item If $vz-wy\ne 0$, then $N''=1$. 
    \item If $vz-wy=0$ and $z(u'-u)=w(x-x')$, then $N''=p$.

\item If $vz-wy=0$ and $z(u'-u)\ne w(x-x')$, then $N''=0$.
\end{itemize}
\end{lemma}


So, combining Propositions \ref{propo2.1}, \ref{propo2.2}, and Lemma \ref{lm23} gives
\[I\le \frac{|P|^2}{p^3}+p^{3-\epsilon}|P|.\]
In other words,
\[I-II-III-IV\ll \frac{1}{p^6}\cdot p^{3-\epsilon}|P|,\]
and 
 \[\left\vert I(P, X)- \frac{|P||X|}{p^3} \right\vert\le p^{\frac{3-\epsilon}{2}}|P|^{\frac{1}{2}}|X|^{\frac{1}{2}}.\]
 This completes the proof.
\end{proof}

The rest of this section is devoted to prove Propositions \ref{propo2.1}, \ref{propo2.2}, and Lemma \ref{lm23}.

We first start with Proposition \ref{propo2.2} and Lemma \ref{lm23}, since they are straightforward. 

\begin{proof}[Proof of Proposition \ref{propo2.2}]
    By fixing $(x, y, z)\in A$ and $(u', v', w')\in B$, then $y'$ and $v$ are determined uniquely by 
    \[yw'+zv'=y'w+vz', ~vz-wy=0, ~zu'+xw'-z'u-x'w=0.\]
With each $u\in \mathbb{F}_p$, $x'$ is determined uniquely. Thus, $N'\le p|A||B|$.
\end{proof}
\begin{proof}[Proof of Lemma \ref{lm23}]
A direct computation shows that 
\[b=\frac{v'-v}{w}=\frac{y'-y}{z}.\]
Note that $z(u'-u)+w(x-x')=a(vz-wy)$. If $vz-wy\ne 0$, then 
\[a=\frac{z(u'-u)+w(x-x')}{vz-wy}.\]

For such $a$, $c$ is determined uniquely. 

If $vz-wy=0$ and $z(u'-u)=w(x-x')$, there are $p$ such matrices $\theta$. 

If $vz-wy=0$ and $z(u'-u)\ne w(x-x')$, there is no such matrix $\theta$. 
\end{proof}

We now move to the proof of Proposition \ref{propo2.1}.

By repeating the proof of \cite[Theorem 2.1]{Hetal}, the following weighted version can be obtained. 
\begin{lemma}\label{lm3.3}
    Let $U, V$ be two sets in $\mathbb{F}_p^4$, and $F\colon U\to \mathbb{N}$, $G\colon V\to \mathbb{N}$. 
   Define 
   \[M=\sum_{\substack{u\in U, v\in V\\ u\cdot v=0}}F(u)G(v).\]Then we have 
    \begin{align*}
    &\left\vert M-\frac{(\sum_{u\in U}F(u))(\sum_{v\in V}G(v))}{p}\right\vert\le  p^2\left(\sum_{u\in U}|F(u)|^2\right)^{\frac{1}{2}}\cdot \left(\sum_{v\in V}|G(v)|^2\right)^{\frac{1}{2}}.     
    \end{align*}
\end{lemma}

In our setting, we partition the set $A$ into $A_\lambda$ and the set $B$ into $B_\lambda$, for $\lambda\in \mathbb{F}_p$, such that any two points in $A_\lambda$ have the same last coordinate which is equal to $\lambda$, and the same applies to $B_\lambda$. 

Recall $\pi_{23}\colon \mathbb{F}_p^3\to \mathbb{F}_p^2$ be the projection onto the two last coordinates. 

For $A_\lambda\subset \mathbb{F}_p^3$, we define $\overline{A_\lambda}=\pi_{23}(A_\lambda)$ . For each $x\in \overline{A_\lambda}\subset \mathbb{F}_p^2$, define $f(x)=\pi_{23}^{-1}(x)\cap A_\lambda$.

%by $f(x)$, we mean the number of points $y$ in $A_\lambda$ such that $\pi_{23}(y)=x$. 

For $B_\lambda\subset \mathbb{F}_p^3$, we define $\overline{B_\lambda}=\pi_{23}(B_\lambda)$ . For each $x\in \overline{B_\lambda}$, define $g(x)=\pi_{23}^{-1}(x)\cap B_\lambda$.
%by $g(x)$, we mean the number of points $y$ in $B_\lambda$ such that $\pi_{23}(y)=x$. 

Let $N_{\lambda, \beta}$ be the number of quadruples $(a, a', b, b')\in \overline{A_\lambda}\times \overline{B_\lambda}\times\overline{A_\beta}\times \overline{B_\beta}$ such that $a\cdot b'=a'\cdot b$, counted with multiplicity, i.e. 
\[N_{\lambda, \beta}=\sum_{a, b, a', b'\colon a\cdot b'=a'\cdot b}f(a)g(a')f(b)g(b')\]

We want to bound $N_{\lambda, \beta}$ from above. The next result will be proved by using the above lemma.
\begin{lemma}\label{weight-H}
    For fixed $\lambda, \beta$, we have 
    \begin{align*}
         &N_{\lambda, \beta}\le \frac{1}{p}\l(\sum_{x\in \overline{A_\lambda}}f(x)\r) \l(\sum_{x\in \overline{A_\beta}}f(x)\r) \l(\sum_{y\in \overline{B_\lambda}}g(y)\r) \l(\sum_{y\in \overline{B_\beta}}g(y)\r)\\
         &+p\l(\sum_{x\in \overline{A_\lambda}}|f(x)|^{2}\r)^{\frac{1}{2}}\l(\sum_{x\in \overline{A_\beta}}|f(x)|^2\r)^{\frac{1}{2}}\l(\sum_{y\in \overline{B_\lambda}}|g(y)|^2\r)^{\frac{1}{2}}\l(\sum_{y\in \overline{B_\beta}}|g(y)|^2\r)^{\frac{1}{2}}.
    \end{align*}
\end{lemma}
\begin{proof}
To prove this lemma, we first enlarge the set $A_\lambda, A_\beta, B_\lambda$, and $B_\beta$. More precisely, define 
\[U=\{t(a, a')\colon a\in \overline{A_{\lambda}}, ~a'\in \overline{B_{\lambda}},~ t\ne 0\},\]
and
\[V=\{t(b, b')\colon b\in \overline{A_{\beta}}, ~b'\in \overline{B_{\beta}},~ t\ne 0\}.\]
Also define 
\[F(t(a, a'))=f(a)g(a'), ~~G(t(b, b'))=f(b)g(b')\]
for all $t(a, a')\in U$ and $t(b, b')\in V$.
%We need to keep in mind that for each $x\in \overline{A_\lambda}$, it has been assigned with the multiplicity $f(x)$, and each $y\in \overline{B_\beta}$, it has been assigned with the multiplicity $g(y)$. By the definition, we mean $A'$ and $B'$ are multi-set, and for each $t\ne 0$ and $x\in \overline{A_\lambda}$, the multiplicity of $tx$ is $f(tx)=f(x)$. The same applies for $B'$ with the function $g$. 

If $(a, a', b, b')\in \overline{A_\lambda}\times \overline{B_\lambda}\times\overline{A_\beta}\times \overline{B_\beta}$ such that $a\cdot b'=a'\cdot b$, then 
\[(a, a')\cdot (-b', b)=0,\]
which implies
\[t(a, a')\cdot t'(-b', b)=0\]
for all $t, t'\ne 0$.

Therefore, 
\[N_{\lambda, \beta}=\frac{1}{(p-1)^2}M,\]
where 
\[M=\sum_{u\in U, v\in V, u\cdot v=0}F(u)G(v).\]
Applying Lemma \ref{lm3.3} gives us 
\begin{align*}
         &N_{\lambda, \beta}\le \frac{1}{p}\l(\sum_{x\in \overline{A_\lambda}}f(x)\r) \l(\sum_{x\in \overline{A_\beta}}f(x)\r) \l(\sum_{y\in \overline{B_\lambda}}g(y)\r) \l(\sum_{y\in \overline{B_\beta}}g(y)\r)\\
         &+p\l(\sum_{x\in \overline{A_\lambda}}|f(x)|^{2}\r)^{\frac{1}{2}}\l(\sum_{x\in \overline{A_\beta}}|f(x)|^2\r)^{\frac{1}{2}}\l(\sum_{y\in \overline{B_\lambda}}|g(y)|^2\r)^{\frac{1}{2}}\l(\sum_{y\in \overline{B_\beta}}|g(y)|^2\r)^{\frac{1}{2}}.
    \end{align*}
This completes the proof. 
\end{proof}
Using this lemma, we take the sum over all pairs $(\lambda, \beta)$, and obtain 
\[N\le \frac{|A|^2|B|^2}{p}+p \left(\sum_{x\in \mathbb{F}_p^2}f(x)^2\right)\cdot \left(\sum_{y \in \mathbb{F}_p^2}g(y)^2\right),\]
by using the Cauchy--Schwarz inequality. 

The trivial upper bound that $f(x), g(y)\le p$ gives 
\[N\le \frac{|A|^2|B|^2}{p}+p^3|A||B|.\]
If we assume in addition that either $\max_{x}f(x)\le p^{1-\epsilon}$ or $\max_{x}g(x)\le p^{1-\epsilon}$, then 
\[N\le \frac{|A|^2|B|^2}{p}+p^{3-\epsilon}|A||B|.\]
This completes the proof. $\square$







\section{Proof of Theorem \ref{main-H}}
With Theorem \ref{H-inci} in hand, one can prove Theorem \ref{main-H} easily. To see this, we set $B=E$ and $A=X(E)$. We now estimate the number of incidences between $A\times B$ and $X$. It is clear that $I(A\times B, X)=|X||B|$. Applying Theorem \ref{H-inci} on the number of incidences between $A\times B$ and $X$, we obtain 
\[I(A\times B, X)\le \frac{|A||B||X|}{p^3}+p^{\frac{3-\epsilon}{2}}|X|^{\frac{1}{2}}|A|^{\frac{1}{2}}|B|^{\frac{1}{2}}.\]
Putting the lower and upper bounds together and solving the inequality give us the desired bound. 
%\section{Open questions}
%\paragraph{Open questions:} We end this section with some open questions. 

%{\bf 1.}
% Based on Theorems \ref{thm2} and \ref{thm3'}, it is plausible to make the following conjecture in the continuous. 
%\begin{conjecture}
%Let $E\subset \mathbb{R}^2$ and $S\subset SL_2(\mathbb{R})$ be compact sets with the corresponding natural metrics. We have 
%\[\dim_H(S(E))\ge \min \left\lbrace 2, ~\dim_H(E)+\dim_H(E)-2 \right\rbrace.\]
%In particular, if $\dim_H(S)+\dim_H(E)>4$, then the set $S(E)$ has positive Lebesgue measure. 
%\end{conjecture}
%\begin{conjecture}
%Let $E\subset \mathbb{R}^2$ and $S\subset SL_2(\mathbb{R})$ be compact sets with the corresponding natural metrics. Assume that $\dim_H(E)>1$ and $\dim_H(S)>2$, then there exists $x\in E$ such that the set $S(E-x)$ has positive Lebesgue measure. 
%\end{conjecture}

%{\bf 2.} In the statement of Theorem \ref{thm3'}, the condition that $S$ is contained entirely in any subgroups of $SL_2(\mathbb{F}_p)$ is needed to guarantee that $\Ebf(S, S)\ll |S|^{3-\epsilon}$. We did not try to compute the precise value of $\epsilon$ from Bourgain and Gamburd's proof, but it should be very small. This leads to two interesting questions of clarifying necessary and sufficient conditions on $S$ such that the inequality $\Ebf(S, S)\ll |S|^{3-\epsilon}$ holds, and finding explicit values for $\epsilon$, possibly with combinatorial arguments and results from incidence geometry. 

%A set $S\subset SL_2(\mathbb{F}_p)$ is called \textit{non-degenerate} if there exists $\epsilon>0$ such that $\Ebf(S, S)\ll |S|^{3-\epsilon}$, and \textit{degenerate} otherwise. Theorem \ref{thm3} and Theorem \ref{thm3'} provide an example of non-degenerate sets $S$. This leads to a question of finding necessary and sufficient conditions for a set to be non-degenerate, which we hope to address in a subsequent paper. 



\section{Acknowledgements}
Thang Pham would like to thank the Vietnam Institute for Advanced Study in Mathematics (VIASM) for the hospitality and for the excellent working condition.

Norbert Hegyv\'{a}ri was supported by the National Research, Development and Innovation Office NKFIH Grant No K-146387. Alex Iosevich was partially supported by the National Science Foundation,  NSF DMS 2154232. Thang Pham was partially supported by ERC Advanced Grant no. 882971, ``GeoScape", and by the Erd\H{o}s Center.
\begin{thebibliography}{}

\bibitem{B.N.P.08} L. Babai, N. Nikolov, and L. Pyber, \textit{Product growth and mixing in finite groups}, Proceedings of the Nineteenth Annual ACM-SIAM Symposium on Discrete Algorithms, (2008), 248--257.

 \bibitem{Besicovitch20} A.~S. Besicovitch, {\em {Sur deux questions d'int{\'e}grabilit{\'e} des fonctions}}, Journal de la Soci{\'e}t{\'e} de Physique et de Math{\'e}matique de l’Universit{\'e} de Perm, {\bf 2} (1919), 105--123.

\bibitem{Besicovitch28}A.~S. Besicovitch, {\em On {K}akeya's problem and a similar one}, Mathematische Zeitschrift, {\bf 27}(1) (1927), 312--320.

 \bibitem{BesicovitchRado68} A.~S. Besicovitch and R.~Rado, {\em A plane set of measure zero containing circumferences of every radius}, Journal of the London Mathematical Society, {\bf 43} (1968), 717--719.


 \bibitem{Bourgain85} J.~Bourgain, {\em Estimations de certaines fonctions maximales}, Comptes rendus de l'Académie des sciences, S\'{e}rie 1, Math\'{e}matique, \textbf{312}(8) (1985), 499--512.

% \bibitem{B.12} J. Bourgain, \textit{A modular Szemerédi–Trotter theorem for hyperbolas}, Comptes Rendus. Mathématique, \textbf{350}(17-18) (2012), 793--796.

\bibitem{BKT1}
 J.~Bourgain, N.~Katz, and T.~Tao, {\it A sum-product estimate in finite fields, and applications}, Geometric \& Functional Analysis, {\bf 14} (2004), 27--57.

\bibitem{B.A.08} J. Bourgain and A. Gamburd, \textit{Uniform expansion bounds for Cayley graphs of $\SL_2(\Fbb_p)$}, Annals of Mathematics, \textbf{167} (2008) 625--642.


\bibitem{covert} D. Covert, D. Hart, A. Iosevich, D. Koh, and M. Rudnev, {\it Generalized incidence theorems, homogeneous forms and sum-product estimates in finite fields}, European Journal of Combinatorics, {\bf 31}(1) (2010), 306–-319.

%\bibitem{dvir}
%Z. Dvir, \textit{On the size of Kakeya sets in finite fields}, Journal of the American Mathematical Society, \textbf{22}(4) (2009), 1093--1097.
\bibitem{Hetal}
D. Hart, A. Iosevich, D. Koh, and M. Rudnev. {\it Averages over hyperplanes, sum-product theory in vector spaces over finite fields and the Erd\H{o}s-Falconer distance conjecture}, Transactions of the American Mathematical Society, \textbf{363}(6) (2011), 3255--3275.

\bibitem{HH0}
N. Hegyvári and F. Hennecart, \textit{A note on Freiman models in Heisenberg groups}, Israel Journal of Mathematics, \textbf{189}(1) (2012), 397--411.
\bibitem{HH1}
N. Hegyv\'{a}ri and F. Hennecart, \textit{A structure result for bricks in Heisenberg groups}, Journal of Number Theory, \textbf{133}(9) (2013), 2999--3006.
\bibitem{HH2}
N. Hegyv\'{a}ri and F. Hennecart, \textit{Expansion for cubes in the Heisenberg group}, Forum Mathematicum, \textbf{30}(1) (2018), 227--236.





% \bibitem{H.I.08} D. Hart and A. Iosevich, {\it Sums and products in finite fields: an integral geometric viewpoint}, Contemporary Mathematics, {\bf 464} (2008).

% \bibitem{Davies72} R.~O. Davies, {\em Another thin set of circles}, J. Lond. Math. Soc., {\bf 5} (1972), 191--192.

% \bibitem{Falconer85book} K.~J. Falconer, {\em The Geometry of Fractal Sets}, Cambridge Tracts in Mathematics, Cambridge University Press, (1985).

% \bibitem{I.R.07} A. Iosevich and M. Rudnev, \textit{Erd\"{o}s distance problem in vector spaces over finite fields}, Transactions of the American Mathematical Society, \textbf{359}(12) (2007), 6127--6142.
\bibitem{IRZ}
A. Iosevich, M. Rudnev, and Y. Zhai, \textit{Areas of triangles and Beck’s theorem in planes over finite fields}, Combinatorica, \textbf{35}(3) (2015), 295--308.

\bibitem{Iosevichetal} A. Iosevich, P. Mattila, E. Palsson, M. Q. Pham, T. Pham, S. Senger, and C. Y. Shen, \textit{Packing sets in euclidean space by affine transformations}, \href{https://arxiv.org/pdf/2405.03087}{arXiv:2405.03087} (2024).

% \bibitem{Keleti17} T.~Keleti, {\em Small union with large set of centers}, in Recent Developments in Fractals and Related Fields, J.~Barral and S.~Seuret, eds., Cham, (2017), Springer International Publishing, 189--206.

\bibitem{Kinney68} J.~R. Kinney, {\em A thin set of circles}, American Mathematical Monthly, {\bf 75} (1968), 1077--1081.

\bibitem{L.P.V.18} B. Lund, T. Pham, and L. A. Vinh, \textit{Distinct spreads in vector spaces over finite fields}, Discrete Applied Mathematics, \textbf{239} (2018), 154--158.

\bibitem{LuPe} B. Lund and G. Petridis, \textit{Bisectors and Pinned Distances}, Discrete \& Computational Geometry, \textbf{64} (2020), 995--1012. 



% \bibitem{L.N.97} R. Lidl and H. Niederreiter, \textit{Finite fields}, Cambridge University Press, (1997).

\bibitem{Marstrand87} J.~M. Marstrand, {\em {P}acking circles in the plane}, Proceedings of the London Mathematical Society. Third Series, {\bf 55} (1987), 37--58.

% \bibitem{Mattila15} P.~Mattila, {\em {F}ourier {A}nalysis and {H}ausdorff {D}imension}, Cambridge Studies in Advanced Mathematics, Cambridge University Press, 2015.

\bibitem{M.T.04}  G. Mockenhaupt and T. Tao. {\it Restriction and Kakeya phenomena for finite fields}, Duke Mathematical Journal, {\bf 121}(1) (2004), 35--74.

\bibitem{M.P.13} G. L. Mullen and D. Panario, \textit{Handbook of Finite Fields}, Chapman and Hall/CRC, (2013).


%\bibitem{O.16} R. Oberlin, \textit{Unions of lines in $F^n$}, Mathematika, \textbf{62} (2016), 738--752.


% \bibitem{P.T.V.19} N. D. Phuong, T. Pham and L. A. Vinh, \textit{Incidences between planes over finite fields}, Proceedings of the American Mathematical Society, \textbf{147} (2019), 2185--2196.

% \bibitem{P.X.23} T. Pham and B. Xue, \textit{New-type Quasirandom Groups and Applications}, \href{https://arxiv.org/abs/2308.01504}{arXiv:2308.01504} (2023).

% \bibitem{P.Y.23} T. Pham and S. Yoo, \textit{Intersection patterns and incidence theorems}, \href{https://arxiv.org/abs/2304.08004}{arXiv:2304.08004} (2023).

\bibitem{P.V.17} T. Pham and L. A. Vinh, \textit{Some combinatorial number theory problems over finite valuation rings}, Illinois Journal of Mathematics, \textbf{61}(1-2) (2017), 243--257.

\bibitem{Pham2024}  T.~Pham, {\it Triangles with one fixed side–length, a Furstenberg-type problem, and incidences in finite vector spaces}, Forum Mathematicum, \textbf{37}(2) (2025), 359--372.

\bibitem{PS}
T. Pham and S. Yoo, \textit{Intersection patterns, incidence theorems, and distance problems}, arXiv:2304.08004 (2025).
 \bibitem{R.14} M. Rudnev, \textit{On the number of incidences between points and planes in three dimensions}, Combinatorica, \textbf{38} (2018), 219--254.


 \bibitem{SdZ} S. Stevens and F. Zeeuw, \textit{An improved point-line incidence bound over arbitrary fields}, Bulletin of the London Mathematical Society, \textbf{49}(5) (2017), 842--858.

\bibitem{Talagrand80} M.~Talagrand, {\em Sur la mesure de la projection d'un compact set certaines families de cercles}, Bulletin des Sciences Math{\'e}matiques, {\bf 104} (1980), 225--231.
 
\bibitem{lavinh} L. A. Vinh. {\it On the volume set of point sets in vector spaces over finite fields}, Proceedings of the American Mathematical Society, {\bf 141}(90) (2013), 3067--3071.
\bibitem{tao}
T. Tao, \textit{Product set estimates for non-commutative groups}, Combinatorica, \textbf{28}(5) (2008), 547--594.
\bibitem{Wolff97} T.~Wolff, {\em A {K}akeya-type problem for circles}, American Journal of Mathematics, \textbf{119}(5) (1997), 985--1026.

\bibitem{Wolff96} T.~Wolff, {\em Recent work connected with the Kakeya problem}, American Mathematical Society, (1999), 129--162.

%\bibitem{Z.23}  J.~Zahl, {\em Union of lines in $\Rbb^n$}, Mathematika, {\bf 69}(2) (2023), 473--481.

\end{thebibliography}

\end{document}
